% !TeX program = xelatex
\documentclass[UTF8,11pt,oneside,openany,fontset=windows]{ctexbook}
\usepackage[a4paper,top=22mm,bottom=22mm,left=22mm,right=20mm,headheight=15pt]{geometry}
\usepackage{amsmath,amssymb,mathtools,bm,booktabs,longtable,array,enumitem,needspace}
\usepackage{xcolor,fancyhdr,hyperref}
\definecolor{ink}{HTML}{163B60}
\definecolor{accent}{HTML}{226A88}
\hypersetup{colorlinks=true,linkcolor=ink,urlcolor=accent,bookmarksnumbered=true,pdftitle={全国大学生数学竞赛：微积分公式定理与真题详解},pdfauthor={学习资料整理}}
\ctexset{chapter={format=\huge\bfseries\color{ink},beforeskip=14pt,afterskip=24pt},section={format=\Large\bfseries\color{ink}},subsection={format=\large\bfseries}}
\setlength{\parskip}{4pt}
\linespread{1.17}
\setlist{nosep,leftmargin=2em,topsep=4pt}
\setcounter{tocdepth}{2}
\allowdisplaybreaks[2]
\emergencystretch=2em
\pagestyle{fancy}\fancyhf{}
\fancyhead[L]{\small 全国大学生数学竞赛 · 非数学类初赛}
\fancyhead[R]{\small\nouppercase{\leftmark}}
\fancyfoot[C]{\thepage}
\renewcommand{\headrulewidth}{0.3pt}
\newcommand{\dd}{\,\mathrm d}
\newcommand{\ee}{\mathrm e}
\newcommand{\ii}{\mathrm i}
\newcommand{\R}{\mathbb R}
\newcommand{\N}{\mathbb N}
\newcommand{\grad}{\nabla}
\DeclareMathOperator{\sgn}{sgn}
\DeclareMathOperator{\curl}{curl}
\DeclareMathOperator{\arsinh}{arsinh}
\DeclareMathOperator{\arcosh}{arcosh}
\newcommand{\src}[1]{\hyperref[src:#1]{[S#1]}}
\newcommand{\srcentry}[1]{\phantomsection\label{src:#1}[S#1]}
\newcounter{problem}
\newcommand{\problem}[3]{\par\bigskip\Needspace{4\baselineskip}\refstepcounter{problem}\subsection*{真题 #1\quad #2}\addcontentsline{toc}{subsection}{真题 #1 #2}\label{q:#1}\noindent{\small\color{accent}来源：#3}\par}
\newcommand{\method}[1]{\par\noindent\textbf{选法：}#1\par}
\newcommand{\sol}{\par\noindent\textbf{详解：}}
\newcommand{\trap}[1]{\par\noindent\textbf{易错与迁移：}#1\par}
\newcommand{\ans}[1]{\boxed{#1}}
\newcommand{\ruleitem}[2]{\par\noindent\textbf{#1}\quad #2\par}
\begin{document}
\begin{titlepage}
\thispagestyle{empty}
\vspace*{20mm}
{\color{accent}\Large 公式要记得住，题目要做得出\par}
\vspace{10mm}
{\Huge\bfseries\color{ink}全国大学生数学竞赛\par}
\vspace{5mm}
{\huge\bfseries 微积分公式、定理\par 与历年真题详解\par}
\vspace{12mm}
{\Large 非数学类 A/B 初赛复习用\par}
\vspace{9mm}
\hrule\vspace{6mm}
\begin{tabular}{@{}rl@{}}
教材主线 & 上海交通大学数学科学学院微积分课程组编\\[3pt]
实际附件 & 《大学数学·微积分》上册第三版、下册第二版\\[3pt]
内容组织 & 公式与条件 → 识别方法 → 真题 → 独立推导\\[3pt]
题目规模 & 84 道（组）已核对来源的历年题，含 A/B 共题及补赛题\\[3pt]
真题年代 & 选自 2009—2025 年间；各题注明届次、类别与题号\\[3pt]
整理日期 & 2026 年 10 月 31 日
\end{tabular}
\vfill
{\small 本资料按竞赛考点重新组织，不是教材逐页摘录，也不是竞赛主办方出版物。题干作简明转述，解答独立编写；完整原卷入口见附录。A/B 卷重复题只计一次，多问同一题按一组计。}
\end{titlepage}
\frontmatter
\chapter*{使用说明}
\addcontentsline{toc}{chapter}{使用说明}
这份资料围绕你提供的两册教材，覆盖非数学类初赛常用的微积分、向量与空间几何、常微分方程。先读每个公式的\textbf{适用条件}，再遮住解答做对应真题。看懂答案后，要能复述“为什么选择这个方法”和“哪一步用到了条件”。

附件核对结果：文件名为“26高数”的 PDF 实际为\textbf{上册第三版}，另一附件为\textbf{下册第二版}。下表采用书内印刷页码，不能直接当作 PDF 页码。不同版本章节细节会有差别，故按主题对照。上册正文第 1 页对应附件 PDF 第 13 页；下册对应 PDF 第 7 页。

\begin{longtable}{p{.23\textwidth}p{.34\textwidth}p{.33\textwidth}}
\toprule 本资料 & 教材对应 & 复习重点\\\midrule\endhead
第 1 章 & 上册第 1、2 章，1—100 页 & 极限、连续、等价量、数列递推\\
第 2、3 章 & 上册第 3、4 章，101—220 页 & 求导、泰勒、中值定理、证明\\
第 4 章 & 上册第 5 章，221—316 页 & 积分、反常积分、几何应用\\
第 5 章 & 上册第 6 章，317 页起 & 一阶方程、二阶常系数方程\\
第 6 章 & 下册第 7 章，1—54 页 & 向量、直线、平面、二次曲面\\
第 7 章 & 下册第 8 章，55—119 页 & 偏导、隐函数、极值、方向导数\\
第 8 章 & 下册第 9 章，120—178 页 & 重积分、换元、对称性\\
第 9 章 & 下册第 10 章，179—237 页 & 线面积分、格林、高斯、斯托克斯\\
第 10 章 & 下册第 11 章，238 页起 & 数项、幂级数、傅里叶级数\\\bottomrule
\end{longtable}

\textbf{三层使用法。}第一遍建立公式框架，先做计算型填空题；第二遍完整书写计算大题和证明题；第三遍回到真题目录，以 150 分钟为一个训练单元混合抽题。150 分钟、100 分见所核对的 2025 年原卷；你参赛当年的具体安排仍以当年通知为准。

\textbf{A/B 使用范围。}两类均可用本资料复习基础与综合题。文中不把某一届的试题差别直接当作永久考纲；较难的存在性证明、傅里叶唯一性、积分与级数交换等标为提高内容。若学校发布了当年大纲，以该大纲确定取舍。

\textbf{符号。}所有角度用弧度；$C^k$ 表示到 $k$ 阶导数连续；$\ln$ 为自然对数；$\arctan$ 取主值；$\dd s,\dd S$ 为弧长、面积元，均无方向；$\dd x,\dd y,\dd z$ 与有向线面积分的符号须随定向处理。数列下标和积分变量仅在各题内部使用。

\textbf{核对原则。}真题来源编号如 \src{17A} 指向附录。早期校方资料库及近年转载卷中，个别参考答案存在符号或排版错误，本资料以原题条件重新推导。没有把自行编造的例子标成真题；个别用于说明定理的短例会明确写“示范”或“反例”。
\tableofcontents
\mainmatter
\chapter{极限、连续与渐近分析}
\section{函数与极限的基本工具}
\ruleitem{基础整理}{复合前检查定义域：$\ln u$ 要求 $u>0$，偶次根号要求被开方数非负，分母非零；反函数的定义域是原函数的值域。奇偶函数的定义域须关于原点对称。奇函数在对称区间可积时积分为零，偶函数积分为半区间的两倍。}
\[
\sin 2x=2\sin x\cos x,\quad 1-\cos x=2\sin^2(x/2),\quad
\sin^2x=\frac{1-\cos2x}{2},\quad \cos^2x=\frac{1+\cos2x}{2}.
\]
\[
\sinh x=\frac{\ee^x-\ee^{-x}}2,\quad\cosh x=\frac{\ee^x+\ee^{-x}}2,\quad
\cosh^2x-\sinh^2x=1.
\]
\ruleitem{极限定义}{数列 $a_n\to A$：$\forall\varepsilon>0$，$\exists N$，$n>N$ 时 $|a_n-A|<\varepsilon$。函数 $f(x)\to A (x\to x_0)$：将“$n>N$”换成“$0<|x-x_0|<\delta$”。极限只管去心邻域，通常不决定 $f(x_0)$。}
\ruleitem{等价刻画}{海涅定理：$f(x)\to A$ 当且仅当对每个趋于 $x_0$、且不等于 $x_0$ 的定义域内数列 $x_n$，均有 $f(x_n)\to A$。找两条数列得到不同极限，可以否定极限存在。柯西准则：实数列收敛当且仅当其尾部项彼此任意接近。}
\ruleitem{运算与保号}{有限极限可作四则运算；商的分母极限须非零。极限为正的量最终为正。无穷大不是实数，$0\cdot\infty$、$\infty-\infty$ 等必须先改写。}
\ruleitem{夹逼与单调有界}{若 $a_n\le b_n\le c_n$ 且外侧同趋于 $L$，则 $b_n\to L$。实数单调有界数列必收敛；递推题先证明“有界、单调”，再把极限代入连续递推式。}
\ruleitem{Stolz 定理（常用型）}{若 $b_n$ 严格递增且趋于 $+\infty$，$\frac{a_{n+1}-a_n}{b_{n+1}-b_n}\to L$，则 $a_n/b_n\to L$，允许 $L$ 为扩展实数。不可省略分母条件。}
\ruleitem{黎曼和}{若 $f$ 在 $[a,b]$ 黎曼可积，均匀分割有
$\displaystyle \frac{b-a}{n}\sum_{k=1}^n f\!\left(a+\frac{k(b-a)}n\right)\to\int_a^b f(x)\dd x$。识别“$1/n$、$k/n$、求和上限 $n$”。}

\section{等价无穷小与泰勒阶数}
记 $f=o(g)$ 表示 $f/g\to0$，$f=O(g)$ 表示比值在所考察邻域有界，$f\sim g$ 表示比值趋于 1（分母在去心邻域非零）。当 $x\to0$：
\[
\sin x\sim\tan x\sim\arcsin x\sim\arctan x\sim x,\quad
\ln(1+x)\sim\ee^x-1\sim x,
\]
\[
1-\cos x\sim\frac{x^2}2,\quad (1+x)^\alpha-1\sim\alpha x (\alpha\ne0),\quad
 a^x-1\sim x\ln a (a>0,a\ne1).
\]
加减式必须先检查抵消。$\sin x\sim x$ 并不允许把 $\sin x-x$ 直接当作零。需要时展开到第一个非零项：
\begin{align*}
\ee^x&=1+x+\frac{x^2}2+\frac{x^3}6+\frac{x^4}{24}+o(x^4),\\
\ln(1+x)&=x-\frac{x^2}2+\frac{x^3}3-\frac{x^4}4+o(x^4),\\
\sin x&=x-\frac{x^3}6+\frac{x^5}{120}+o(x^5),\\
\cos x&=1-\frac{x^2}2+\frac{x^4}{24}-\frac{x^6}{720}+o(x^6),\\
\tan x&=x+\frac{x^3}3+\frac{2x^5}{15}+o(x^5),\\
\arctan x&=x-\frac{x^3}3+\frac{x^5}5+o(x^5),\\
(1+x)^\alpha&=1+\alpha x+\frac{\alpha(\alpha-1)}2x^2+\frac{\alpha(\alpha-1)(\alpha-2)}6x^3+o(x^3).
\end{align*}
\ruleitem{幂指极限}{底数最终为正时 $u(x)^{v(x)}=\exp(v(x)\ln u(x))$。若 $h\to0$、$vh\to L$，则 $(1+h)^v\to\ee^L$。这里的 $1^\infty$ 只是型态，不是答案。}
\ruleitem{增长阶}{当 $x\to+\infty$，固定 $a,b>0$，$\ln x=o(x^a)$，$x^a=o(\ee^{bx})$。含 $n!$ 可先取对数、夹逼或用递推比值，不必一律使用斯特林公式。}
\ruleitem{积分上限小量}{若 $f(t)\sim c t^m$、$m>-1$，且 $u\to0^+$，则 $\int_0^u f(t)\dd t\sim c u^{m+1}/(m+1)$（$c\ne0$）。可以由比较估计证明，不是将积分号随意“移出极限”。}

\section{连续性与闭区间定理}
连续：$\lim_{x\to a}f(x)=f(a)$，须同时检查定义、极限、相等。第一类间断点左右极限都有限；左右不等是跳跃，相等但与函数值不符或未定义是可去；其余为第二类。

闭区间连续函数有界、取得最大最小值，且一致连续；介值定理保证取遍两端函数值之间的数。零点存在常用 $f(a)f(b)<0$。仅有开区间连续不够，如 $1/x$ 在 $(0,1)$ 无界。复合函数的极限定理在外层连续时可直接使用。

\section{数列、幂指与黎曼和真题}
\problem{01}{连乘递推的望远镜化}{2010 年第 2 届，非数学类，一(1)，\src{02}}
设 $|a|<1$，求 $P_n=\prod_{k=0}^n(1+a^{2^k})$ 的极限。
\method{看到指数 $1,2,4,8,\ldots$，使用平方差连续相消。}
\sol $a=0$ 时结论直接成立。一般地，
\[(1-a)P_n=(1-a)(1+a)(1+a^2)\cdots(1+a^{2^n})=1-a^{2^{n+1}}.\]
由于 $a^{2^{n+1}}\to0$，故 $\ans{\lim P_n=1/(1-a)}$。迁移：形如 $1+x^{2^k}$ 的积先补上 $1-x$，比取对数求和更直接。

\problem{02}{三角连乘}{2011 年第 3 届，非数学类，一(2)，\src{03}}
求 $\lim_{n\to\infty}\prod_{k=1}^n\cos(\theta/2^k)$。
\sol 多次使用倍角公式，
\[2^n\sin(\theta/2^n)\prod_{k=1}^n\cos(\theta/2^k)=\sin\theta.\]
固定 $\theta\ne0$ 时，充分大的 $n$ 使分母非零，且 $2^n\sin(\theta/2^n)\to\theta$。答案为 $\ans{\sin\theta/\theta}$；$\theta=0$ 时各因子均为 1，极限为 1。

\problem{03}{阶乘的弱指数}{2012 年第 4 届，非数学类，一(1)，\src{04}}
求 $\lim_{n\to\infty}(n!)^{1/n^2}$。
\sol 取对数，$0\le \ln(n!)/n^2\le n\ln n/n^2=\ln n/n\to0$。再取指数得 $\ans{1}$。不要只因 $n!$ 增长快就把原式判断为无穷大；外部指数同样关键。

\problem{04}{带整数周期的幂指极限}{2013 年第 5 届，非数学类，一(1)，\src{05}}
求 $\lim_{n\to\infty}[1+\sin(\pi\sqrt{4n^2+1})]^n$。
\sol 去掉整数周期，令
\[\delta_n=\pi(\sqrt{4n^2+1}-2n)=\frac\pi{\sqrt{4n^2+1}+2n},\quad n\delta_n\to\frac\pi4.\]
于是 $n\ln(1+\sin\delta_n)\sim n\delta_n\to\pi/4$，底数最终为正，答案 $\ans{\ee^{\pi/4}}$。

\problem{05}{先夹逼，再识别黎曼和}{2015 年第 7 届，非数学类，一(1)，\src{07}}
求 $\displaystyle\lim_{n\to\infty}n\sum_{k=1}^n\frac{\sin(k\pi/n)}{n^2+k}$。
\sol 因正弦非负，
\[\frac1{n+1}\sum_{k=1}^n\sin\frac{k\pi}n
\le n\sum_{k=1}^n\frac{\sin(k\pi/n)}{n^2+k}
\le\frac1n\sum_{k=1}^n\sin\frac{k\pi}n.\]
右端趋于 $\int_0^1\sin(\pi t)\dd t=2/\pi$，左端多一个趋于 1 的因子 $n/(n+1)$，故答案 $\ans{2/\pi}$。注意原题分母是 $n^2+k$，没有根号。

\problem{06}{导数定义进入指数}{2016 年第 8 届，非数学类，一(1)，\src{08}}
$f$ 在 $a$ 可导，$f(a)\ne0$。求 $\lim[f(a+1/n)/f(a)]^n$。
\sol 由可导性，$f(a+1/n)/f(a)=1+f'(a)/(nf(a))+o(1/n)$，故答案
$\ans{\exp(f'(a)/f(a))}$。即使 $f(a)<0$，此比值也最终趋近正数 1，实数幂表达仍可用取对数法。

\problem{07}{先抽出导数定义}{2016 年第 8 届，非数学类，一(2)，\src{08}}
已知 $f(1)=0$、$f'(1)$ 存在，求
\[\lim_{x\to0}\frac{f(\sin^2x+\cos x)\tan3x}{(\ee^{x^2}-1)\sin x}.\]
\sol $u=\sin^2x+\cos x-1=\tfrac12x^2+o(x^2)$，可导性给出 $f(1+u)=f'(1)u+o(u)$，其余因子分别为 $3x+o(x)$、$x^2+o(x^2)$、$x+o(x)$。所以极限为 $\ans{3f'(1)/2}$。此写法也涵盖 $f'(1)=0$，不应错误写成 $f(1+u)\sim f'(1)u$。

\problem{08}{差分中的中值定理}{2018 年第 10 届，非数学类，一(1)，\src{10}}
固定 $0<\alpha<1$，求 $\lim[(n+1)^\alpha-n^\alpha]$。
\sol 由拉格朗日中值定理，差为 $\alpha\xi_n^{\alpha-1}$，其中 $n<\xi_n<n+1$，故答案为 $\ans0$。也可写成 $n^\alpha[(1+1/n)^\alpha-1]\sim\alpha n^{\alpha-1}$。

\problem{09}{倒数与阶乘共同消去递推}{2020 年第 12 届，非数学类，二，\src{12}}
$a_1=1$，$a_{n+1}=a_n/[(n+1)(a_n+1)]$。求 $\lim n!a_n$。
\sol 归纳知 $a_n>0$。令 $b_n=1/(n!a_n)$，则
\[b_{n+1}=\frac{a_n+1}{n!a_n}=b_n+\frac1{n!},\qquad b_1=1.\]
所以 $b_n=\sum_{k=0}^{n-1}1/k!\to\ee$，从而 $\ans{n!a_n\to1/\ee}$。原递推含 $a_n+1$ 时，倒数常把非线性变成差分。

\problem{10}{递推趋零时如何求速度}{第 13 届初赛补赛，非数学类，一(1)，\src{13}}
$x_0=1$，$x_n=\ln(1+x_{n-1})$，求 $\lim nx_n$。
\sol 对 $t>0$ 有 $0<\ln(1+t)<t$，故 $x_n$ 正、递减，极限满足 $L=\ln(1+L)$，只能为 0。由泰勒式，
\[\frac1{x_{n+1}}-\frac1{x_n}=\frac{x_n-\ln(1+x_n)}{x_n\ln(1+x_n)}\to\frac12.\]
取 $b_n=n$ 使用 Stolz，得 $(1/x_n)/n\to1/2$，故 $\ans{nx_n\to2}$。只有 $x_n\to0$ 还不足以给出这个速度。

\section{近年卷：泰勒展开的抵消与组合}
\problem{11}{无穷远的标准幂指型}{2023 年第 15 届，非数学 B 类，一(1)，\src{15B}}
求 $\displaystyle\lim_{x\to+\infty}\left(\frac{x+3}{x+2}\right)^{2x-1}$。
\sol 对数为 $(2x-1)\ln(1+1/(x+2))\to2$，故答案 $\ans{\ee^2}$。用 $1+u$ 形式比机械求导更清楚。

\problem{12}{先抵消二阶项}{2025 年第 17 届，非数学 A 类，一(1)，\src{17A}}
设 $f(x)=x-\ln(1+x)-\tfrac12x\sin x$，$g(x)=ax^3$，且 $\lim_{x\to0}f/g=1$，求 $a$。
\sol 展开至三阶：$x-\ln(1+x)=x^2/2-x^3/3+o(x^3)$，$x\sin x/2=x^2/2+o(x^3)$，故 $f=-x^3/3+o(x^3)$。于是 $(-1/3)/a=1$，$\ans{a=-1/3}$。

\problem{13}{变上限积分的阶数}{2025 年第 17 届，非数学 A 类，一(3)，\src{17A}}
求 $\displaystyle\lim_{x\to0}\frac{\int_0^{\sin^2x}\ln(1+t)\dd t}{((1+x^3)^{1/3}-1)\sin x}$。
\sol 上限 $u=\sin^2x\sim x^2$；因 $\ln(1+t)\sim t$，分子 $\sim u^2/2\sim x^4/2$。分母 $\sim(x^3/3)x=x^4/3$，得 $\ans{3/2}$。先在积分变量 $t$ 上展开，再代入上限，避免把 $t$ 与 $x$ 混为一谈。

\problem{14}{复合展开与首项抵消}{2024 年第 16 届，非数学 B 类，一(1)，\src{16B}}
求 $\displaystyle\lim_{x\to0}\frac{\ln(1+x^2)-x^2}{\sin^4x}$。
\sol 令 $u=x^2$，$\ln(1+u)-u=-u^2/2+o(u^2)$，分子为 $-x^4/2+o(x^4)$，分母 $\sim x^4$，故答案 $\ans{-1/2}$。这是“先看抵消，再选展开阶数”的最小模板。

\chapter{导数、微分与高阶展开}
\section{基本公式与求导结构}
导数是线性近似的系数：$f(a+h)=f(a)+f'(a)h+o(h)$。可导必连续，连续未必可导，如 $|x|$ 在 0。左右导数相等且有限才可导。微分 $\dd y=f'(x)\dd x$；复合函数的一阶微分形式不变，高阶微分一般不能直接照搬。
\begin{longtable}{p{.44\textwidth}p{.45\textwidth}}
\toprule 函数与导数 & 条件或补充\\\midrule\endhead
$(x^\alpha)'=\alpha x^{\alpha-1}$ & 一般实指数取 $x>0$，整数次幂可按其定义域扩展\\
$(\ee^x)'=\ee^x$；$(a^x)'=a^x\ln a$ & $a>0$\\
$(\ln|x|)'=1/x$；$(\log_a x)'=1/(x\ln a)$ & 前者 $x\ne0$；后者 $x>0,a>0,a\ne1$\\
$(\sin x)'=\cos x$；$(\cos x)'=-\sin x$ & $x\in\R$\\
$(\tan x)'=\sec^2x$；$(\cot x)'=-\csc^2x$ & 避开原函数无定义处\\
$(\sec x)'=\sec x\tan x$ & $(\csc x)'=-\csc x\cot x$\\
$(\arcsin x)'=1/\sqrt{1-x^2}$ & $(\arccos x)'=-1/\sqrt{1-x^2}$，$|x|<1$\\
$(\arctan x)'=1/(1+x^2)$ & $(\operatorname{arccot}x)'=-1/(1+x^2)$，采用常见主值约定\\
$(\sinh x)'=\cosh x$ & $(\cosh x)'=\sinh x$；$(\tanh x)'=1/\cosh^2x$\\
$(\arsinh x)'=1/\sqrt{1+x^2}$ & $\arsinh x=\ln(x+\sqrt{1+x^2})$\\\bottomrule
\end{longtable}
\ruleitem{线性、乘积、商与复合}{
$(uv)'=u'v+uv'$；$(u/v)'=(u'v-uv')/v^2$（$v\ne0$）；$(f\circ g)'=f'(g)g'$。
乘积与复合交错出现时，先标清最外层运算。}
\ruleitem{对数求导}{若 $y=u(x)^{v(x)}$ 且 $u>0$，则 $y'=u^v[v'\ln u+vu'/u]$。多因子乘除也可先求 $\ln|y|$，再还原 $y'$。}
\ruleitem{反函数}{若 $y=f(x)$ 在某区间可逆且 $f'(x)\ne0$，则 $(f^{-1})'(y)=1/f'(x)$，二阶为 $-f''(x)/(f'(x))^3$。每次求导要说明自变量是谁。}
\ruleitem{参数方程}{若 $x=x(t),y=y(t)$ 且 $x'(t)\ne0$，
\[
\frac{\dd y}{\dd x}=\frac{y'}{x'},\qquad
\frac{\dd^2y}{\dd x^2}=\frac{x'y''-y'x''}{(x')^3}.
\]}
\ruleitem{一元隐函数}{对 $F(x,y)=0$，若 $F\in C^1$ 且 $F_y\ne0$，局部存在 $y=y(x)$，
\[
y'=-\frac{F_x}{F_y},\qquad
y''=-\frac{F_{xx}+2F_{xy}y'+F_{yy}(y')^2}{F_y}\quad(F\in C^2).
\]}
\ruleitem{高阶乘积公式}{若相关导数存在，
$\displaystyle(uv)^{(n)}=\sum_{k=0}^n\binom nk u^{(k)}v^{(n-k)}$。
常用 $\frac{\dd^n}{\dd x^n}\ee^{ax}=a^n\ee^{ax}$，$\frac{\dd^n}{\dd x^n}\sin(ax+b)=a^n\sin(ax+b+n\pi/2)$。}
\ruleitem{系数提取}{若 $f(a+h)=\sum_{k=0}^n c_kh^k+o(h^n)$ 且相应可导，则 $f^{(n)}(a)=n!c_n$。不是直接取 $c_n$。}

\section{泰勒公式与几何应用}
\ruleitem{Peano 余项}{在 $a$ 具有 $n$ 阶导数时，
$\displaystyle f(a+h)=\sum_{k=0}^n\frac{f^{(k)}(a)}{k!}h^k+o(h^n)$。
这是局部极限工具，不能据此直接估算远处误差。}
\ruleitem{Lagrange 余项}{若 $f$ 在 $a,x$ 之间有到 $n+1$ 阶的充分光滑性（例如 $C^{n+1}$），则
\[
f(x)=\sum_{k=0}^n\frac{f^{(k)}(a)}{k!}(x-a)^k+
\frac{f^{(n+1)}(\xi)}{(n+1)!}(x-a)^{n+1},
\]
$\xi$ 在 $a,x$ 之间；若 $|f^{(n+1)}|\le M$，误差绝对值不超过 $M|x-a|^{n+1}/(n+1)!$。}
\ruleitem{积分余项}{若 $f\in C^{n+1}$，余项也可写为
$\displaystyle R_n(x)=\frac1{n!}\int_a^x f^{(n+1)}(t)(x-t)^n\dd t$。
带权积分不等式常比不定的 $\xi$ 更方便。}
\ruleitem{切线、法线}{曲线 $y=f(x)$ 在 $a$ 的切线 $y-f(a)=f'(a)(x-a)$；法线用 $(x-a)+f'(a)(y-f(a))=0$，此式也适用于 $f'(a)=0$。参数曲线可用方向向量 $(x',y')$ 写切线，避免遗漏竖直切线。}
\ruleitem{弧长与曲率}{
\[
\dd s=\sqrt{1+(y')^2}\dd x,quad
\kappa=\frac{|y''|}{(1+(y')^2)^{3/2}},\quad
\kappa=\frac{|x'y''-y'x''|}{((x')^2+(y')^2)^{3/2}}.
\]
后式要求曲线正则，即 $(x',y')\ne(0,0)$。曲率半径为 $1/\kappa$（$\kappa\ne0$）。}
\ruleitem{数值求根（补充）}{牛顿迭代 $x_{n+1}=x_n-f(x_n)/f'(x_n)$。若 $f\in C^2$、$f(r)=0$、$f'(r)\ne0$，初值足够靠近简单根 $r$，则局部二阶收敛；远离根时可能发散，不能无条件使用。}

\section{高阶导数与隐函数真题}
\problem{15}{指数与正弦乘积的四阶导数}{2016 年第 8 届，非数学类，一(4)，\src{08}}
求 $f(x)=\ee^x\sin2x$ 的 $f^{(4)}(0)$。
\method{只取 $x^4$ 的系数，比连续四次求导省算。}
\sol $\ee^x=1+x+x^2/2+x^3/6+O(x^4)$，$\sin2x=2x-4x^3/3+O(x^5)$。
乘积的 $x^4$ 系数为 $2/6-4/3=-1$，故 $\ans{f^{(4)}(0)=-24}$。
也可用复数 $(1+2\ii)^4=-7-24\ii$ 复核虚部。部分网络参考解把符号写为正，这里以计算结果为准。

\problem{16}{零点重数与莱布尼茨公式}{2020 年第 12 届，非数学类，一(2)，\src{12}}
求 $f(x)=(x+1)^n\ee^{-x^2}$ 的 $f^{(n)}(-1)$（$n\ge1$）。
\sol 对 $(x+1)^n$ 求不足 $n$ 次导数，代入 $-1$ 均为 0。莱布尼茨和中只剩该因子求 $n$ 次、指数因子不求导的一项，故 $\ans{n!/\ee}$。等价地令 $h=x+1$，$f(-1+h)=\ee^{-1}h^n+O(h^{n+1})$。

\problem{17}{部分分式给出任意阶导数}{2023 年第 15 届，非数学 A 类，一(3)，\src{15A}}
设 $f(x)=1/(x^2-3x+2)$，求 $f^{(n)}(0)$。
\sol $f(x)=1/(1-x)-1/(2-x)$，在 $|x|<1$ 展开得
\[f(x)=\sum_{n=0}^\infty\left(1-\frac1{2^{n+1}}\right)x^n.\]
故 $\ans{f^{(n)}(0)=n!(1-2^{-n-1})}$。先分解两个简单极点，避免直接对有理函数反复求导。

\problem{18}{含变上限积分的二阶隐导数}{2024 年第 16 届，非数学 B 类，一(3)，\src{16B}}
$\ee^{3y}+\int_0^{x+y}\cos(t^2)\dd t=1$ 确定过原点的隐函数，求 $y''(0)$。
\sol 一次求导：
\[3\ee^{3y}y'+\cos((x+y)^2)(1+y')=0.\]
代入原点得 $4y'(0)+1=0$。再求导时，$\cos((x+y)^2)$ 的导数为
$-2(x+y)\sin((x+y)^2)(1+y')$，原点值为 0，所以
\[9(y'(0))^2+4y''(0)=0,\qquad \ans{y''(0)=-9/64}.\]
切勿漏掉 $\ee^{3y}$ 求导得到的 $(y')^2$ 项。

\problem{19}{参数曲线与隐函数合用}{2018 年第 10 届，非数学类，一(2)，\src{10}}
$x=t+\cos t$，$\ee^y+ty+\sin t=1$，求 $t=0$ 时的切线。
\sol $t=0$ 时 $(x,y)=(1,0)$，且 $x'(0)=1$。第二式对 $t$ 求导：
$(\ee^y+t)y'+y+\cos t=0$，故 $y'(0)=-1$。切线为 $\ans{x+y-1=0}$。

\problem{20}{小量恒等式同时确定函数值与导数}{第 14 届初赛补赛二，一(2)，\src{14}}
在 1 附近可导的 $f$ 满足
$f(1+x)-3f(1-x)=4+2x+o(x)$，求曲线在 $x=1$ 处切线。
\sol 左端展开为 $-2f(1)+4f'(1)x+o(x)$。比较常数、一次项得 $f(1)=-2$、$f'(1)=1/2$，所以 $y+2=(x-1)/2$，即 $\ans{x-2y-5=0}$。
“比较系数”在此可严格化：先令 $x\to0$，再减去常数项、除以 $x$ 求极限。

\problem{21}{零点重数让复合求导变简单}{2025 年第 17 届，非数学 B 类，一(1)，\src{17B}}
$y=x^2\ee^{\sin x}$，求 $y''(0)$。
\sol $\ee^{\sin x}=1+O(x)$，所以 $y=x^2+O(x^3)$；函数光滑，由二阶系数得 $\ans{y''(0)=2}$。若用莱布尼茨公式，也仅有 $(x^2)''\ee^{\sin0}$ 一项保留。

\chapter{中值定理、证明题与不等式}
\section{定理条件与使用信号}
\ruleitem{费马定理}{若 $x_0$ 是定义区间内部的局部极值点且 $f'(x_0)$ 存在，则 $f'(x_0)=0$。驻点未必是极值，$x^3$ 在 0 即为反例；端点极值不要求导数为零。}
\ruleitem{Rolle}{若 $f\in C[a,b]$、在 $(a,b)$ 可导，且 $f(a)=f(b)$，则存在 $\xi\in(a,b)$ 使 $f'(\xi)=0$。有 $m+1$ 个不同零点，反复使用可得到相应导数的零点。}
\ruleitem{Lagrange}{同样的连续、可导条件下，存在 $\xi\in(a,b)$：$f(b)-f(a)=f'(\xi)(b-a)$。用来把“函数差”变成“导数界”。}
\ruleitem{Cauchy}{若 $f,g$ 均满足上述条件，存在 $\xi$：
$[f(b)-f(a)]g'(\xi)=[g(b)-g(a)]f'(\xi)$。
要写比值型还需相应分母非零，常用 $g'(x)\ne0$ 保证。}
\ruleitem{L'Hôpital}{当分子、分母同趋于 0，或分母绝对值趋于无穷的相应型，函数在去心单侧邻域可导、$g'\ne0$，且 $f'/g'$ 有有限或无穷极限时，可推出 $f/g$ 的该极限。$0\cdot\infty$、$\infty-\infty$、幂指型须先改写。导数比极限不存在，不代表原比值不存在。}

\textbf{证明题的四个入口。}
目标有 $f'$ 的零点，考虑 Rolle；有 $f'+pf$，考虑乘积分因子 $\ee^{\int p}$；有两个不同点的导数，先找中间分点再在两侧分别用中值定理；有积分与函数值，先把积分变为原函数。

\section{单调、凹凸、零点与基本不等式}
$f'\ge0$ 可推出不减，$f'>0$ 推出严格增；反向不成立，例如 $x^3$ 严格增却在 0 导数为 0。二阶导数 $f''\ge0$ 时称为凸函数，本文采用弦在图像上方的约定：
\[
f(\lambda x+(1-\lambda)y)\le\lambda f(x)+(1-\lambda)f(y),\quad0\le\lambda\le1.
\]
可微凸函数位于每条切线之上；凹函数不等号反向。拐点要求凹凸性改变，$f''=0$ 仅是候选条件。闭区间求最值需比较驻点、不可导点及端点。
\begin{align*}
\ee^x&\ge1+x &&(x\in\R),\\
\ln(1+x)&\le x &&(x>-1),\\
|\sin x|&\le |x| &&(x\in\R),\\
2|ab|&\le \lambda a^2+\lambda^{-1}b^2 &&(\lambda>0),\\
\left|\sum a_kb_k\right|^2&\le\left(\sum a_k^2\right)\left(\sum b_k^2\right),\\
\left|\int_a^b uv\dd x\right|^2&\le\int_a^b u^2\dd x\int_a^b v^2\dd x.
\end{align*}
Cauchy 不等式的等号要求两列成比例；积分版本对连续函数要求处处成比例（一般可积情形为几乎处处）。Jensen 不等式是凸性向有限加权和的推广；使用前检查权重非负、和为 1 及凸性区间。

\ruleitem{根的数量}{存在用介值；唯一用严格单调；最多 $m$ 个常用反证：若有 $m+1$ 个不同根，经 $m$ 次 Rolle 迫使 $f^{(m)}$ 有零点。约束导数的符号有助于排除多根。}
\ruleitem{高阶导数判别}{若 $f'(a)=\cdots=f^{(m-1)}(a)=0$，$f^{(m)}(a)\ne0$，则首个非零阶 $m$ 为偶数时按该导数正负判断极小、极大；奇数时不是极值。需有相应 Taylor 展开。}

\section{零点、辅助函数与积分不等式真题}
\problem{22}{严格凸函数恰有两个零点}{2010 年第 2 届，非数学类，二，\src{02}}
设 $f''(x)>0$，$\lim_{x\to+\infty}f'(x)=\alpha>0$，$\lim_{x\to-\infty}f'(x)=\beta<0$，且某点 $f(x_0)<0$。证明恰有两个实根。
\sol 充分大的 $x$ 有 $f'(x)>\alpha/2$，积分或中值定理给出 $f(x)\to+\infty$。充分小的 $x$ 有 $f'(x)<\beta/2<0$；向负方向积分同样得到 $f(x)\to+\infty$。与 $f(x_0)<0$ 配合，左右两侧各有一根。
若有三个不同根，两次 Rolle 使 $f''$ 在某点为 0，与严格正矛盾，故恰有两根。
\trap{负半轴的导数为负，向左移动时函数反而升高。不要把 $x\to-\infty$ 时的符号判反。}

\problem{23}{端点条件迫使三阶导数取指定值}{2011 年第 3 届，非数学类，三，\src{03}}
$f\in C^3[-1,1]$，$f(-1)=0$，$f(1)=1$，$f'(0)=0$。证明存在 $\xi\in(-1,1)$，$f'''(\xi)=3$。
\sol 在 0 作二阶 Taylor 展开，余项取三阶：
\[1=f(0)+\tfrac12f''(0)+\tfrac16f'''(\eta_+),\qquad
0=f(0)+\tfrac12f''(0)-\tfrac16f'''(\eta_-),\]
其中 $\eta_-\in(-1,0)$，$\eta_+\in(0,1)$。两式相减得 $f'''(\eta_+)+f'''(\eta_-)=6$。连续函数 $f'''$ 在两值之间必取平均数 3（若两值已等于 3，直接取该点）。
\trap{两次 Taylor 余项的中间点一般不同，不能擅自合并为同一个点。这里靠连续介值完成最后一步。}

\problem{24}{两个不同点的倒数关系}{2016 年第 8 届，非数学类，五，\src{08}}
连续函数 $f$ 满足 $I=\int_0^1 f\dd x\ne0$。证明存在不同的 $x_1,x_2\in(0,1)$，使 $1/f(x_1)+1/f(x_2)=2/I$。
\method{先找“累计积分的一半”，再在两段分别用中值定理。}
\sol 令 $F(x)=I^{-1}\int_0^x f(t)\dd t$，$F(0)=0,F(1)=1$。介值定理给出 $c\in(0,1)$、$F(c)=1/2$。在 $[0,c]$ 与 $[c,1]$ 分别用 Lagrange，得到
\[f(x_1)=\frac I{2c},\quad f(x_2)=\frac I{2(1-c)},\quad 0<x_1<c<x_2<1.\]
两值均非零，倒数相加即为 $2/I$。$f$ 不必恒正，$F$ 不必单调；介值已经足够。

\problem{25}{先介值，再两次中值}{2020 年第 12 届，非数学类，三，\src{12}}
$f\in C[0,1]$，在 $(0,1)$ 可导，$f(0)=0,f(1)=1$。证明存在 $c\in(0,1)$ 使 $f(c)=2-3c$；并存在不同 $\xi,\eta\in(0,1)$，使 $(1+f'(\xi))(1+f'(\eta))=4$。
\sol 对 $g(x)=f(x)+3x-2$，有 $g(0)=-2,g(1)=2$，故存在该 $c$。分别在两侧使用 Lagrange：
\[f'(\xi)=\frac{2-3c}{c},\qquad f'(\eta)=\frac{3c-1}{1-c},\quad \xi<c<\eta.\]
于是 $1+f'(\xi)=2(1-c)/c$，$1+f'(\eta)=2c/(1-c)$，乘积为 4。

\problem{26}{由零积分到非线性导数等式}{2023 年第 15 届，非数学 B 类，五，\src{15B}}
$f$ 在 $[0,1]$ 可导，$f(0)>0,f(1)>0,\int_0^1 f=0$。证明至少有两个零点，且某 $\xi\in(0,1)$ 满足 $f'(\xi)+3f^3(\xi)=0$。
\sol 若 $f\ge0$，由端点附近连续且为正，积分必严格大于零。因此有 $x_0\in(0,1)$ 使 $f(x_0)<0$。两侧各由介值定理得零点 $a<b$。
构造 $G(x)=f(x)\exp(\int_0^x3f^2(t)\dd t)$，则 $G(a)=G(b)=0$。Rolle 给出 $G'(\xi)=0$，而
\[G'(x)=\exp\left(\int_0^x3f^2(t)\dd t\right)[f'(x)+3f^3(x)],\]
指数因子非零，所证即得。看到 $f'+3f^3=f'+(3f^2)f$，就能识别积分因子。

\problem{27}{分部积分加 Cauchy，连等号也要证明}{2023 年第 15 届，非数学 A 类五／B 类六，共题，\src{15A}\src{15B}}
$f\in C^1[0,1]$，$f(0)=0$。证明 $\int_0^1 f^2\le4\int_0^1(1-x)^2(f')^2$，并求等号条件。
\sol 记两积分为 $I,J$。分部积分得
\[I=[(x-1)f^2]_0^1+2\int_0^1(1-x)ff'\dd x
=2\int_0^1(1-x)ff'\dd x\le2\sqrt{IJ}.\]
$I=0$ 时 $f\equiv0$；$I>0$ 时平方除以 $I$ 得 $I\le4J$。
若取等且 $I>0$，Cauchy 等号要求 $(1-x)f'=cf$。代回 $I=2\int(1-x)ff'$ 得 $c=1/2$。于是 $f=C(1-x)^{-1/2}$（$x<1$），由 $f(0)=0$ 得 $C=0$，与 $I>0$ 矛盾。所以等号当且仅当 $\ans{f\equiv0}$。

\problem{28}{导数满足 Lipschitz 条件的全局估计}{2024 年第 16 届，非数学 A 类五／B 类六，共题，\src{16A}\src{16B}}
$f\in C^1(\R)$，$f\ge0$，且 $|f'(x)-f'(y)|\le M|x-y|$，$M>0$。证明 $(f'(x))^2\le2Mf(x)$。
\sol 由微积分基本定理，对任意实数 $h$，
\[|f(x+h)-f(x)-f'(x)h|\le\frac M2h^2.\]
故 $0\le f(x+h)\le f(x)+f'(x)h+Mh^2/2$。右端二次式对每个 $h$ 均非负，取 $h=-f'(x)/M$，得 $0\le f(x)-(f'(x))^2/(2M)$。
\trap{定义域是整个实轴，所以可自由取最优 $h$。若只在有限区间非负，不能直接照搬此选择。}

\problem{29}{把目标不等式变成两个单调函数}{2016 年第 8 届，非数学类，二，\src{08}}
$f$ 在 $[0,1]$ 可导，$f(0)=0$，$0<f'(x)<1$。证明对 $0<a<1$，$(\int_0^a f)^2>\int_0^a f^3$。
\sol 对 $x>0$，$f(x)>0$。令 $A(x)=\int_0^x f(t)\dd t$，
$G(x)=2A(x)-f^2(x)$。因 $G(0)=0$、$G'=2f(1-f')>0$，有 $G(x)>0$。
再令 $H(x)=A^2(x)-\int_0^x f^3(t)\dd t$，则 $H(0)=0$，
\[H'(x)=2A(x)f(x)-f^3(x)=f(x)G(x)>0.\]
故 $H(a)>0$。目标里同时出现“积分的平方”和“三次方积分”，先整体求导，再对剩余括号求导，是可迁移的构造路径。

\chapter{一元积分、反常积分与应用}
\section{积分定理与必备原函数}
\ruleitem{黎曼积分}{有界函数的积分通过分割、取样、积分和的极限定义。连续函数、只有有限个间断点的有界函数均可积；单调有界函数可积。可积不必连续。若 $f\le g$ 则 $\int f\le\int g$；$|\int f|\le\int|f|$。}
\ruleitem{积分中值定理}{若 $f$ 连续、$g$ 可积且不变号，存在 $\xi\in[a,b]$ 使 $\int_a^b fg=f(\xi)\int_a^b g$。特别地 $\int_a^b f=f(\xi)(b-a)$。一般只保证闭区间内的点。}
\ruleitem{微积分基本定理}{若 $f$ 连续，$F(x)=\int_a^x f(t)\dd t$，则 $F'=f$。若 $\Phi'=f$，则 $\int_a^b f=\Phi(b)-\Phi(a)$。变上、下限：$\frac{\dd}{\dd x}\int_{u(x)}^{v(x)}f(t)\dd t=f(v)v'-f(u)u'$。}
\ruleitem{换元与分部}{
\[
\int f(\varphi(t))\varphi'(t)\dd t=\int f(x)\dd x,\qquad
\int u\dd v=uv-\int v\dd u.
\]
定积分换元必须同时换上下限；分部积分优先让对数、反三角、多项式求导降阶。}
以下积分常数均为 $C$；在各自定义域的连通区间内使用。参数 $a>0$：
\begin{align*}
\int x^p\dd x&=\frac{x^{p+1}}{p+1}+C\quad(p\ne-1),&\int\frac{\dd x}x&=\ln|x|+C,\\
\int\ee^{ax}\dd x&=\frac{\ee^{ax}}a+C,&\int\ln x\dd x&=x\ln x-x+C,\\
\int\sin x\dd x&=-\cos x+C,&\int\cos x\dd x&=\sin x+C,\\
\int\tan x\dd x&=-\ln|\cos x|+C,&\int\cot x\dd x&=\ln|\sin x|+C,\\
\int\sec^2x\dd x&=\tan x+C,&\int\csc^2x\dd x&=-\cot x+C,\\
\int\sec x\dd x&=\ln|\sec x+\tan x|+C,&\int\csc x\dd x&=\ln|\csc x-\cot x|+C,\\
\int\frac{\dd x}{a^2+x^2}&=\frac1a\arctan\frac xa+C,&
\int\frac{\dd x}{\sqrt{a^2-x^2}}&=\arcsin\frac xa+C,\\
\int\frac{\dd x}{x^2-a^2}&=\frac1{2a}\ln\left|\frac{x-a}{x+a}\right|+C,\\
\int\frac{\dd x}{\sqrt{x^2+a^2}}&=\ln(x+\sqrt{x^2+a^2})+C,\\
\int\frac{\dd x}{\sqrt{x^2-a^2}}&=\ln|x+\sqrt{x^2-a^2}|+C,\\
\int\ee^{ax}\cos bx\dd x&=\frac{\ee^{ax}(a\cos bx+b\sin bx)}{a^2+b^2}+C,\\
\int\ee^{ax}\sin bx\dd x&=\frac{\ee^{ax}(a\sin bx-b\cos bx)}{a^2+b^2}+C.
\end{align*}
最后两式要求 $a^2+b^2\ne0$。常见根式换元：$\sqrt{a^2-x^2}$ 用 $x=a\sin t$；$\sqrt{a^2+x^2}$ 用 $x=a\tan t$ 或 $a\sinh t$；$\sqrt{x^2-a^2}$ 可用 $x=a\sec t$。选择区间使根号及雅可比符号明确。三角有理式可用 $t=\tan(x/2)$：
\[\sin x=\frac{2t}{1+t^2},\quad\cos x=\frac{1-t^2}{1+t^2},\quad\dd x=\frac{2\dd t}{1+t^2}.\]
有理函数先多项式除法，再按实一次因子与不可约二次因子分式分解。

\section{对称性、周期与反常收敛}
\ruleitem{区间反射}{对可积 $f$，$\int_a^b f(x)\dd x=\int_a^b f(a+b-x)\dd x$。两式相加常把复杂函数消去。若 $g(a+b-x)=g(x)$，则 $\int_a^b xg(x)\dd x=\frac{a+b}2\int_a^b g(x)\dd x$。}
\ruleitem{周期与递推}{周期为 $T$ 的可积函数，任一完整周期积分相同。令 $I_n=\int_0^{\pi/2}\sin^nx\dd x$，有 $I_n=(n-1)I_{n-2}/n$，$I_0=\pi/2,I_1=1$；余弦幂同值。}
\ruleitem{反常积分定义}{无穷区间或无界被积函数处取单侧极限。若有多个瑕点，必须分别收敛；不能用两个发散部分相消。$\int_{-\infty}^{\infty}f$ 要求两端分别收敛，主值存在不等于反常积分收敛。}
\[
\int_1^\infty x^{-p}\dd x\text{ 收敛}\iff p>1,\qquad
\int_0^1 x^{-p}\dd x\text{ 收敛}\iff p<1.
\]
$\int_0^1|\ln x|^q\dd x$ 对 $q\ge0$ 收敛。比较判别用于最终非负函数；若 $f/g\to c\in(0,\infty)$，则同敛散。$\int|f|<\infty$ 推出 $\int f$ 收敛，反向未必。
\ruleitem{Dirichlet 与 Abel}{若 $\int_a^A f$ 关于 $A$ 有界，$g$ 单调趋于 0，则 $\int_a^\infty fg$ 收敛；这是积分版 Dirichlet。若 $\int f$ 收敛、$g$ 单调有界，则乘积积分收敛（Abel）。振荡函数乘衰减量时先想它们。}
\ruleitem{常用特殊积分（提高）}{
\[
\int_0^\infty\ee^{-ax^2}\dd x=\frac{\sqrt\pi}{2\sqrt a}\ (a>0),\quad
\Gamma(s)=\int_0^\infty x^{s-1}\ee^{-x}\dd x\ (s>0),
\]
$\Gamma(s+1)=s\Gamma(s)$，$\Gamma(1)=1$，$\Gamma(1/2)=\sqrt\pi$；
$B(p,q)=\int_0^1 t^{p-1}(1-t)^{q-1}\dd t=\Gamma(p)\Gamma(q)/\Gamma(p+q)$，$p,q>0$。可识别结构，但证明题不得只报特殊函数名称代替推导。}

\section{几何与物理应用公式}
平面面积 $A=\int_a^b|f-g|\dd x$；极坐标扇形面积 $A=\frac12\int_\alpha^\beta r^2\dd\theta$，要求区域被恰当描绘一次。旋转体圆盘/垫片法 $V=\pi\int(R^2-r^2)\dd x$；柱壳法 $V=2\pi\int(\text{半径})(\text{高度})\dd x$，须避免重复计体积。
\[
L=\int\sqrt{(x')^2+(y')^2}\dd t,\quad
L_{\rm polar}=\int\sqrt{r^2+(r')^2}\dd\theta,\quad
S_{\rm rev}=2\pi\int(\text{到旋转轴的距离})\dd s.
\]
质量 $m=\int\rho\dd s$（线密度），质心为坐标的一阶矩除以质量；转动惯量 $I=\int r_\perp^2\dd m$；变力作功 $W=\int F(x)\dd x$。先写微元，再选积分上下限，比背特例可靠。

\section{积分方程、分部与对称真题}
\problem{30}{把定积分当成一个常数}{2009 年第 1 届，非数学类，一(2)，\src{01}}
连续函数满足 $f(x)=3x^2-2-\int_0^2f(t)\dd t$，求 $f$。
\sol 设 $A=\int_0^2f$，则 $f=3x^2-2-A$。再积分得 $A=8-4-2A$，故 $A=4/3$，$\ans{f(x)=3x^2-10/3}$。此处积分变量 $t$ 是哑变量，整个积分不依赖外部 $x$。

\problem{31}{一个极限常数与一个积分常数}{2025 年第 17 届，非数学 A 类，二，\src{17A}}
连续函数在 0 可导，满足 $f(x)=x^3+x^2\lim_{t\to0}f(t)/t-x\int_0^1f(t)\dd t$，求 $f$。
\sol 记极限为 $a$、积分为 $b$，则 $f=x^3+ax^2-bx$，这也给出 $f(0)=0$。除以 $x$ 取极限得 $a=-b$。再积分：$b=1/4+a/3-b/2$，解得 $b=3/22,a=-3/22$，故
\[\ans{f(x)=x^3-\frac3{22}x^2-\frac3{22}x}.\]
将结果代回两个常数的定义即可复核，而不是仅满足其中一条。

\problem{32}{指数衰减与阶乘递推}{2010 年第 2 届，非数学类，一(3)，\src{02}}
$s>0$、整数 $n\ge1$，求 $I_n=\int_0^\infty\ee^{-sx}x^n\dd x$。
\sol 指数衰减保证收敛，且 $x^n\ee^{-sx}\to0$。先在 $[0,A]$ 分部再令 $A\to\infty$，得 $I_n=(n/s)I_{n-1}$，$I_0=1/s$。所以 $\ans{I_n=n!/s^{n+1}}$。无穷远边界项不能不检查就丢掉。

\problem{33}{绝对值周期乘指数}{2012 年第 4 届，非数学类，二，\src{04}}
求 $I=\int_0^\infty\ee^{-2x}|\sin x|\dd x$。
\sol 按长度 $\pi$ 分段，令 $x=k\pi+t$，每段是第一段的 $\ee^{-2k\pi}$ 倍，且非负可逐段相加。
\[
\int_0^\pi\ee^{-2t}\sin t\dd t=\frac{1+\ee^{-2\pi}}5,\qquad
I=\ans{\frac{1+\ee^{-2\pi}}{5(1-\ee^{-2\pi})}}.
\]
先消绝对值，再利用周期；$\ee^{-2x}|\sin x|$ 本身不是周期函数。

\problem{34}{两次反射消去反三角函数与因子 x}{2013 年第 5 届，非数学类，二，\src{05}}
求 $I=\int_{-\pi}^{\pi}\frac{x\sin x\arctan(\ee^x)}{1+\cos^2x}\dd x$。
\sol $w(x)=x\sin x/(1+\cos^2x)$ 为偶函数，且 $\arctan(\ee^x)+\arctan(\ee^{-x})=\pi/2$。配对 $x,-x$ 得
\[I=\frac\pi2\int_0^\pi\frac{x\sin x}{1+\cos^2x}\dd x.\]
再用 $x\mapsto\pi-x$，$\sin x/(1+\cos^2x)$ 关于 $\pi/2$ 对称，所以
\[I=\frac{\pi^2}4\int_0^\pi\frac{\sin x}{1+\cos^2x}\dd x
=\frac{\pi^2}4[-\arctan(\cos x)]_0^\pi=\ans{\pi^3/8}.\]
第一步利用偶权配对，第二步利用区间中点对称，两种对称不要混用。

\problem{35}{正弦积分为什么不是绝对收敛}{2013 年第 5 届，非数学类，一(2)，\src{05}}
判断 $\int_0^\infty\sin x/x\dd x$ 是否绝对收敛。
\sol 0 处连续延拓后无瑕点。无穷远按 Dirichlet 收敛。对 $k\ge1$，
\[\int_{k\pi}^{(k+1)\pi}\frac{|\sin x|}x\dd x
\ge\frac1{(k+1)\pi}\int_{k\pi}^{(k+1)\pi}|\sin x|\dd x
=\frac2{(k+1)\pi}.\]
对 $k$ 求和发散，所以为\textbf{条件收敛}。不能因 $|\sin x|\le1$、$\int1/x$ 发散，就直接推出原绝对值积分发散；上界发散没有判定力。

\problem{36}{反双曲函数形式的分部积分}{2018 年第 10 届，非数学类，一(3)，\src{10}}
求 $\int\frac{\ln(x+\sqrt{1+x^2})}{(1+x^2)^{3/2}}\dd x$。
\sol 取 $u=\arsinh x$，$\dd v=\dd x/(1+x^2)^{3/2}$，则 $v=x/\sqrt{1+x^2}$、$\dd u=\dd x/\sqrt{1+x^2}$。故结果为
\[\ans{\frac{x\ln(x+\sqrt{1+x^2})}{\sqrt{1+x^2}}-\frac12\ln(1+x^2)+C}.\]
识别 $(x/\sqrt{1+x^2})'=(1+x^2)^{-3/2}$ 是关键。

\problem{37}{对数定积分的降阶}{2024 年第 16 届，非数学 A 类，一(1)，\src{16A}}
求 $\int_0^1\ln(1+x^2)\dd x$。
\sol 分部得 $[x\ln(1+x^2)]_0^1-\int_0^1 2x^2/(1+x^2)\dd x$，将后者写成 $2-2/(1+x^2)$，答案为 $\ans{\ln2-2+\pi/2}$。

\problem{38}{奇函数与半圆面积}{2024 年第 16 届，非数学 B 类，一(2)，\src{16B}}
求 $\int_{-2}^2[x^3\cos^5x+\sqrt{4-x^2}]\dd x$。
\sol 第一项为奇函数，在对称区间积分为 0；第二项为半径 2 的上半圆面积。故结果 $\ans{2\pi}$。无需做三角换元。

\problem{39}{函数关系给出偶部}{2025 年第 17 届，非数学 B 类，一(3)，\src{17B}}
连续函数满足 $f(x)+f(-x)=2$，求 $\int_{-1}^1[f(x)+x^2\ln(x+\sqrt{1+x^2})]\dd x$。
\sol $\arsinh x$ 为奇函数，因此第二项积分为 0。第一项配对后为 $\int_0^1[f(x)+f(-x)]\dd x=2$，答案 $\ans2$。不需要求出 $f$ 的具体形式。

\problem{40}{反常积分先截断再分部}{2025 年第 17 届，非数学 B 类，一(5)，\src{17B}}
求 $\int_1^\infty\frac{\ln x}{(1+x)^2}\dd x$。
\sol 边界项 $-\ln x/(1+x)$ 在无穷远趋于 0，在 1 为 0，故
\[I=\int_1^\infty\frac{\dd x}{x(1+x)}
=\left[\ln\frac{x}{1+x}\right]_1^\infty=\ans{\ln2}.\]
分部后两项 $1/x-1/(1+x)$ 必须作为一个整体积分，不能拆成两个发散积分相减。

\problem{41}{由 Dirichlet 积分推导 sinc 的平方、立方}{2024 年第 16 届，非数学 B 类，五，\src{16B}}
已知 $\int_0^\infty\sin x/x\dd x=\pi/2$，求 $I_2=\int_0^\infty(\sin x/x)^2\dd x$ 与 $I_3=\int_0^\infty(\sin x/x)^3\dd x$。
\sol 先在截断区间分部并检查两端边界，得
\[I_2=\int_0^\infty\frac{\sin2x}{x}\dd x=\ans{\pi/2}.\]
再分部两次：
\begin{align*}
I_3&=\frac32\int_0^\infty\frac{\sin^2x\cos x}{x^2}\dd x\\
&=\frac38\int_0^\infty\frac{\cos x-\cos3x}{x^2}\dd x\\
&=\frac38\int_0^\infty\frac{-\sin x+3\sin3x}{x}\dd x
=\ans{3\pi/8}.
\end{align*}
第一步边界含 $\sin^3x/x^2$，第二步含 $(\cos x-\cos3x)/x$，两端均为 0。最后用换元 $u=3x$；组合式在拆分前应先保证每个所得正弦积分收敛。

\problem{42}{旋转体积最小：先消参数，再求二次函数最值}{2023 年第 15 届，非数学 B 类，二，\src{15B}}
$y=3ax^2+2bx+\ln c$（$c>0$）过原点，在 $[0,1]$ 非负；与横轴及 $x=1$ 围成面积为 1。求绕横轴旋转的最小体积。
\sol 过原点得 $c=1$；面积条件给出 $a+b=1$。于是
\[V=\pi\int_0^1[3ax^2+2(1-a)x]^2\dd x
=\pi\left(\frac{2a^2}{15}+\frac a3+\frac43\right).\]
抛物线最小值在 $a=-5/4$，$b=9/4$，值为 $\ans{9\pi/8}$。此时 $y=(-15x^2+18x)/4\ge0$ 于 $[0,1]$，候选确实满足原约束；若不检验非负性，代数最优未必可行。

\chapter{常微分方程}
\section{一阶方程：先识别，再变形}
\ruleitem{可分离变量}{$y'=p(x)q(y)$：当 $q(y)\ne0$，积分 $\int\dd y/q(y)=\int p(x)\dd x+C$。每个满足 $q(y_0)=0$ 的常值函数须单独检查，除法可能丢解。}
\ruleitem{齐次型}{$y'=F(y/x)$，在 $x\ne0$ 的区间令 $y=vx$，则 $xv'=F(v)-v$，再分离。这里“齐次型”与线性齐次方程是不同概念。}
\ruleitem{一阶线性}{$y'+P(x)y=Q(x)$，连续系数区间内
\[
\mu(x)=\exp\left(\int P(x)\dd x\right),\qquad
y=\mu^{-1}\left(\int\mu Q\dd x+C\right).
\]
公式来自 $(\mu y)'=\mu Q$；积分常数只保留一个独立常数即可。}
\ruleitem{Bernoulli}{$y'+Py=Qy^m$，$m\ne0,1$，在允许除以 $y^m$ 的区间令 $u=y^{1-m}$，得 $u'+(1-m)Pu=(1-m)Q$。检查 $y=0$ 是否为原方程解，非整数幂还需检查实数定义域。}
\ruleitem{恰当微分}{$M(x,y)\dd x+N(x,y)\dd y=0$。在单连通域内 $M,N\in C^1$ 且 $M_y=N_x$，存在势函数 $\Phi$，$\Phi_x=M,\Phi_y=N$，解为 $\Phi=C$。先对一个变量积分，再用另一个偏导确定遗漏的单变量函数。}
\ruleitem{一元积分因子}{若 $(M_y-N_x)/N$ 只依赖 $x$，可取 $\mu(x)=\ee^{\int(M_y-N_x)/N\dd x}$；若 $(N_x-M_y)/M$ 只依赖 $y$，可取相应的 $\mu(y)$。除法处仍需检查非零条件。}
\ruleitem{初值解唯一性}{对 $y'=F(x,y)$，$F$ 连续且关于 $y$ 局部 Lipschitz（例如 $F_y$ 连续），给定初值可保证局部存在唯一解。仅连续通常只保证存在，不能保证唯一。解的区间可能受系数奇点、爆破解限制。}

\section{二阶常系数、降阶与共振}
\ruleitem{线性结构}{$y''+py'+qy=r(x)$ 的通解为 $y_h+y_p$。两个非齐次解之差是齐次解；非齐次解随意相加一般不再是该方程的解。}
特征方程 $\lambda^2+p\lambda+q=0$：
\[
y_h=\begin{cases}
C_1\ee^{\lambda_1x}+C_2\ee^{\lambda_2x},&\lambda_1\ne\lambda_2\text{ 为实根},\\
(C_1+C_2x)\ee^{\lambda x},&\lambda\text{ 为二重实根},\\
\ee^{\alpha x}(C_1\cos\beta x+C_2\sin\beta x),&\alpha\pm\ii\beta,\ \beta\ne0.
\end{cases}
\]
若右端为 $\ee^{\alpha x}P_m(x)$，试取 $y_p=x^k\ee^{\alpha x}Q_m(x)$，$k$ 为 $\alpha$ 在特征方程中的重数。若右端含 $\cos\beta x,\sin\beta x$，两者都要保留，按复根 $\alpha+\ii\beta$ 的重数乘 $x^k$。所谓共振，就是试解与齐次解重合后需要补乘 $x$。
\ruleitem{常用降阶}{不显含 $y$ 的 $y''=F(x,y')$，令 $p=y'$；不显含 $x$ 的 $y''=F(y,y')$，令 $p=p(y)$，$y''=p\dd p/\dd y$，同时检查常值解。}
\ruleitem{Euler 方程}{$x^2y''+axy'+by=r(x)$，在 $x>0$ 令 $t=\ln x$，有 $xy'=Y_t$，$x^2y''=Y_{tt}-Y_t$，化为常系数方程；$x<0$ 可用 $t=\ln|x|$ 分区间处理。}

\section{结构识别真题}
\problem{43}{由三个非齐次解反求方程}{2009 年第 1 届，非数学类，五，\src{01}}
同一二阶常系数非齐次线性方程有三个解：$y_1=x\ee^x+\ee^{2x}$，$y_2=x\ee^x+\ee^{-x}$，$y_3=x\ee^x+\ee^{2x}-\ee^{-x}$。求以最高阶系数为 1 的方程及通解。
\sol $y_1-y_3=\ee^{-x}$，$y_2-y_3=2\ee^{-x}-\ee^{2x}$ 是两个线性无关齐次解，故特征根为 $-1,2$，齐次算子为 $D^2-D-2$。
取特解 $x\ee^x$，代入得 $r(x)=(1-2x)\ee^x$。所以
\[
\ans{y''-y'-2y=(1-2x)\ee^x},\qquad
\ans{y=C_1\ee^{2x}+C_2\ee^{-x}+x\ee^x}.
\]
从解的差得到齐次结构，再找右端，顺序不能颠倒。

\problem{44}{重根对应的两个基本解}{2014 年第 6 届，非数学类，一(1)，\src{06}}
已知 $\ee^x,x\ee^x$ 是某二阶常系数齐次方程的两个解，求方程。
\sol 两解线性无关，对应特征根 1 的重数为 2。因此规范化方程为 $\ans{y''-2y'+y=0}$。只用 $\ee^x$ 不能唯一确定两个特征根；第二个解 $x\ee^x$ 提供重根信息。

\problem{45}{常数右端与初值}{第 14 届初赛补赛二，一(3)，\src{14}}
解 $y''-2y'-3y=1$，$y(0)=0,y'(0)=1$。
\sol 齐次根为 $3,-1$，常数特解为 $-1/3$，故 $y=C_1\ee^{3x}+C_2\ee^{-x}-1/3$。初值给出 $C_1+C_2=1/3$、$3C_1-C_2=1$，解得 $C_1=1/3,C_2=0$，因此 $\ans{y=(\ee^{3x}-1)/3}$。代回原方程和两个初值复核。

\problem{46}{微分形式中识别 y 平方}{2024 年第 16 届，非数学 A 类二／B 类四，共题，\src{16A}\src{16B}}
求 $(x^3-y^2)\dd x+(x^2y+xy)\dd y=0$ 的通解。
\method{只有 $y^2$ 和 $y\dd y$，令 $u=y^2$ 变为一阶线性。}
\sol 以 $x$ 为自变量，$\dd u=2y\dd y$，得
\[x^3-u+\frac{x(x+1)}2u'=0,\qquad
u'-\frac2{x(x+1)}u=-\frac{2x^2}{x+1}.\]
在 $x\ne0,-1$ 的连通区间，积分因子 $\mu=((x+1)/x)^2$，所以
\[(\mu u)'=-2(x+1),\qquad
\ans{\frac{(x+1)^2}{x^2}y^2+(x+1)^2=C}.\]
也可整理为 $(x+1)^2(x^2+y^2)=Cx^2$。后式只是方便表达，并不自动把被除掉的直线恢复为图像解；若讨论微分形式的所有积分曲线，竖直直线 $x=0$、$x=-1$ 也须直接代回原式检查，它们确为积分曲线。过奇点的局部解需按原式分析。

\problem{47}{三角非线性用 sec 代换}{2025 年第 17 届，非数学 A 类，三，\src{17A}}
解 $x\ln x\sin y\,y'+(1-x\cos y)\cos y=0$（实数范围 $x>0$）。
\sol 在 $\cos y\ne0$ 的区间令 $u=\sec y$。除以 $x\cos^2y$ 得
\[\ln x\,u'+\frac ux=1,\qquad (u\ln x)'=1.\]
所以 $\ans{\sec y\ln x=x+C}$，或隐式写成 $\ln x=(x+C)\cos y$。在 $x\ne1$ 的常规区间可按一阶线性求解；写成乘积求导则避免了额外除以 $\ln x$。
另有被 $\cos y$ 除法遗漏的常值解 $\ans{y=\pi/2+k\pi}$（$k\in\mathbb Z$）。隐式式子在 $x=1$ 可能退化，实际延拓必须满足原方程，不能只因代入隐式等式成立便认定为可导解。

\chapter{向量代数与空间解析几何}
\section{向量、距离与夹角公式}
\[
\bm a\cdot\bm b=|\bm a||\bm b|\cos\theta=\sum a_ib_i,\qquad
|\bm a\times\bm b|=|\bm a||\bm b|\sin\theta.
\]
叉积用三阶行列式，方向遵循右手法则；标量三重积 $[\bm a,\bm b,\bm c]=(\bm a\times\bm b)\cdot\bm c$ 是有向体积，取绝对值得平行六面体体积，除以 6 得四面体体积。非零向量垂直等价于点积为零；平行等价于叉积为零。投影长度为 $\bm a\cdot\bm b/|\bm b|$，投影向量为 $(\bm a\cdot\bm b)\bm b/|\bm b|^2$。

\ruleitem{直线与平面}{$\bm r=\bm r_0+t\bm v$；平面 $\bm n\cdot(\bm r-\bm r_0)=0$。方向向量某分量为零时，不应在对称式中除以该分量，应写对应坐标固定。两个平面相交线方向为 $\bm n_1\times\bm n_2$。}
\ruleitem{夹角}{两直线取锐角时 $\cos\theta=|\bm v_1\cdot\bm v_2|/(|\bm v_1||\bm v_2|)$；两平面用法向量同式；线面夹角 $\sin\theta=|\bm v\cdot\bm n|/(|\bm v||\bm n|)$。}
\ruleitem{距离}{
\[
d(P,\Pi)=\frac{|Ax_0+By_0+Cz_0+D|}{\sqrt{A^2+B^2+C^2}},\quad
d(P,l)=\frac{|(\bm P-\bm P_0)\times\bm v|}{|\bm v|}.
\]
两不平行直线间最短距离为
$\displaystyle d=\frac{|(\bm P_2-\bm P_1)\cdot(\bm v_1\times\bm v_2)|}{|\bm v_1\times\bm v_2|}$。
平行时改用点到直线距离，不可把分母为零的式子硬算。}
\ruleitem{平面束}{过两相交平面 $L_1=0,L_2=0$ 的交线的平面为 $\lambda L_1+\mu L_2=0$，$(\lambda,\mu)\ne(0,0)$。写成 $L_1+kL_2=0$ 会遗漏 $L_2=0$ 本身。}

\section{常见曲面与切平面}
\begin{longtable}{p{.45\textwidth}p{.44\textwidth}}
\toprule 标准方程（$a,b,c>0$） & 几何含义\\\midrule\endhead
$x^2/a^2+y^2/b^2+z^2/c^2=1$ & 椭球面\\
$x^2/a^2+y^2/b^2-z^2/c^2=1$ & 单叶双曲面，轴为 $z$ 轴\\
$z^2/c^2-x^2/a^2-y^2/b^2=1$ & 双叶双曲面，轴为 $z$ 轴\\
$z=x^2/a^2+y^2/b^2$ & 椭圆抛物面\\
$z=x^2/a^2-y^2/b^2$ & 双曲抛物面（马鞍面）\\
$x^2/a^2+y^2/b^2-z^2/c^2=0$ & 椭圆锥面\\
缺少某坐标的曲面方程 & 通常为沿该坐标方向延伸的柱面\\\bottomrule
\end{longtable}
通过平移完成平方，把非标准方程化为标准形式；截面法有助于识别内部区域。旋转曲面绕 $z$ 轴时，常把平面曲线的径向变量换为 $\sqrt{x^2+y^2}$。

对光滑曲面 $F(x,y,z)=0$，若 $\grad F(P)\ne0$，切平面为 $\grad F(P)\cdot(\bm r-\bm P)=0$，法线方向为 $\grad F(P)$。图像型 $z=f(x,y)$ 可取法向 $(f_x,f_y,-1)$。两曲面交线的切向为两法向叉积，要求叉积非零。

\section{空间几何真题}
\problem{48}{异面直线距离}{2010 年第 2 届，非数学类，一(5)，\src{02}}
$l_1:x=y,z=0$，$l_2:(x-2)/4=(y-1)/(-2)=(z-3)/(-1)$，求距离。
\sol 取 $P_1=(0,0,0)$、$P_2=(2,1,3)$，$\bm v_1=(1,1,0)$、$\bm v_2=(4,-2,-1)$，有 $\bm v_1\times\bm v_2=(-1,1,-6)$。所以
\[d=\frac{|(2,1,3)\cdot(-1,1,-6)|}{\sqrt{38}}=\ans{\frac{19}{\sqrt{38}}=\frac{\sqrt{38}}2}.\]
取绝对值后距离才非负，不能直接用有向三重积作答案。

\problem{49}{过交线的两个垂直平面}{2012 年第 4 届，非数学类，一(2)，\src{04}}
两平面都过 $2x+y-3z+2=0$ 与 $5x+5y-4z+3=0$ 的交线，且互相垂直，其中一个过 $(4,-3,1)$。求二平面。
\sol 该点代入两原平面左端均为 4，故过该点的平面束参数满足 $\lambda+\mu=0$，得到 $\Pi_1:3x+4y-z+1=0$。
另一个的法向为 $(2\lambda+5\mu,\lambda+5\mu,-3\lambda-4\mu)$，与 $(3,4,-1)$ 点积为 $13\lambda+39\mu=0$，取 $\lambda=-3\mu$，得
\[\ans{\Pi_1:3x+4y-z+1=0,\quad\Pi_2:x-2y-5z+3=0}.\]

\problem{50}{指定平行方向的切平面}{2016 年第 8 届，非数学类，一(5)，\src{08}}
求曲面 $z=x^2/2+y^2$ 上与平面 $2x+2y-z=0$ 平行的切平面。
\sol 法向 $(x,2y,-1)$ 与 $(2,2,-1)$ 平行，因第三分量相同，比例为 1，所以切点 $(2,1,3)$。切平面为 $2(x-2)+2(y-1)-(z-3)=0$，即 $\ans{2x+2y-z=3}$。

\problem{51}{过两点且垂直另一平面}{2025 年第 17 届，非数学 A/B 类，一(2)，共题，\src{17A}\src{17B}}
平面过原点及 $(6,-3,2)$，并垂直于 $4x-y+2z-8=0$，求方程。
\sol 所求法向同时垂直于 $(6,-3,2)$ 与 $(4,-1,2)$，可取它们的叉积 $(-4,-4,6)$，约成 $(2,2,-3)$，所以 $\ans{2x+2y-3z=0}$。常数 $-8$ 影响参照平面位置，不影响它的法向。

\chapter{多元函数微分、隐函数与极值}
\section{极限、偏导与全微分}
\ruleitem{多元极限}{$f(x,y)\to L$ 要求从定义域内任意路径趋近同一点都得到 $L$。两条路径不同即可否定；所有直线路径相同仍不足以证明。证明常用 $r=\sqrt{x^2+y^2}$ 控制 $|f-L|\le\omega(r)\to0$。}
\ruleitem{可微定义}{$f(a+h,b+k)=f(a,b)+Ah+Bk+o(\sqrt{h^2+k^2})$。可微推出连续且偏导存在，$A=f_x,B=f_y$；偏导存在未必连续、更未必可微。若偏导在邻域存在且在该点连续，则可微。}
\ruleitem{链式法则}{若 $z=f(u(x,y),v(x,y))$，各映射可微，则 $z_x=f_u u_x+f_v v_x$，$z_y=f_u u_y+f_v v_y$。二阶求导必须同时求“外层偏导”和“内层系数”的导数。若 $z=f(x(t),y(t))$，$z'=\grad f\cdot\bm r'$。}
\ruleitem{混合偏导}{在邻域内相关二阶偏导连续时，$f_{xy}=f_{yx}$。仅两个混合偏导在一点存在，不足以交换。写 $f_{uv}$ 时 $u,v$ 是外层自变量，不是原来的 $x,y$。}
\ruleitem{全微分}{$\dd z=f_x\dd x+f_y\dd y$；$\dd^2z=f_{xx}\dd x^2+2f_{xy}\dd x\dd y+f_{yy}\dd y^2$，此处把 $x,y$ 视为独立变量。复合路径的二阶导数另有 $f_xx''+f_yy''$ 项，不应遗漏。}
\ruleitem{隐函数}{$F(x,y,z)=0$，$F\in C^1$、$F_z\ne0$，则局部
\[
z_x=-F_x/F_z,\quad z_y=-F_y/F_z,\quad
z_{xy}=-\frac{F_{xy}+F_{xz}z_y+F_{yz}z_x+F_{zz}z_xz_y}{F_z}
\]
（后式要求 $C^2$）。方程组隐函数用关于待解变量的雅可比行列式非零作为条件，再解线性方程组。}

\section{方向导数、泰勒与条件极值}
\ruleitem{方向导数}{沿\textbf{单位向量} $\bm e$，可微时 $D_{\bm e}f=\grad f\cdot\bm e$。最大值 $|\grad f|$，方向为梯度方向（梯度非零时）；最小值为其负值。任意方向导数都存在也不能反推可微。}
\ruleitem{齐次函数 Euler 公式}{若 $f(tx,ty)=t^mf(x,y)$（正的 $t$ 在相应域内），且 $f$ 可微，则 $xf_x+yf_y=mf$。三元同理。先检查缩放性质，不可仅凭式子看起来“次数相同”就使用。}
\ruleitem{二元 Taylor}{在 $C^2$ 条件下，$\bm h=(h,k)^T$，
\[
f(\bm a+\bm h)=f(\bm a)+\grad f(\bm a)\cdot\bm h+\frac12\bm h^TH_f(\bm a)\bm h+o(|\bm h|^2).
\]
二次型或 Hessian 正定/负定决定驻点严格极小/极大；不定则为鞍点。}
驻点处令 $A=f_{xx},B=f_{xy},C=f_{yy}$，$\Delta=AC-B^2$：$\Delta>0,A>0$ 为极小；$\Delta>0,A<0$ 为极大；$\Delta<0$ 为鞍点；$\Delta=0$ \textbf{不能判断}。示范：$x^4+y^4$ 极小，$-x^4-y^4$ 极大，$x^4-y^4$ 鞍点，三者二阶量在原点全为零。

\ruleitem{Lagrange 乘子}{约束 $g=0$ 上求 $f$ 极值，正则点 $\grad g\ne0$ 满足 $\grad f=\lambda\grad g$。多约束要求约束梯度线性无关。它给候选点，须比较函数值；约束奇点与边界也需检查。闭有界区域上连续函数必取最值，可按“内部驻点、边界、角点”分别求。}
\ruleitem{几何}{图像曲面切平面 $z-z_0=f_x(x-x_0)+f_y(y-y_0)$；等值面法向为梯度。曲线 $\bm r(t)$ 的切向为 $\bm r'(t)$，正则点的法平面以该切向作法向。}

\section{二阶链式求导真题}
\problem{52}{径向复合函数的 Laplace 算子}{2010 年第 2 届，非数学类，一(4)，\src{02}}
$r=\sqrt{x^2+y^2}>0$，$g(x,y)=f(1/r)$，$f\in C^2$，求 $g_{xx}+g_{yy}$。
\sol 对径向函数 $h(r)$，$r_x=x/r,r_y=y/r$，故 $g_{xx}+g_{yy}=h''+h'/r$。这里
\[h'=-r^{-2}f'(1/r),\quad h''=2r^{-3}f'(1/r)+r^{-4}f''(1/r).\]
从而 $\ans{g_{xx}+g_{yy}=r^{-4}f''(1/r)+r^{-3}f'(1/r)}$。径向 Laplace 公式取决于空间维数；三维的系数为 $2h'/r$，不能混用。

\problem{53}{指数因子消去一阶项}{2012 年第 4 届，非数学类，一(3)，\src{04}}
令 $z=u(x,y)\ee^{ax+by}$，其中 $u_{xy}=0$。求使任意此类 $u$ 均满足 $z_{xy}-z_x-z_y+z=0$ 的 $a,b$。
\sol 提取指数，左端化为 $u_{xy}+(b-1)u_x+(a-1)u_y+(ab-a-b+1)u$。要对任意 $u$ 成立，需 $a=b=1$，常数项也随之为 0。故 $\ans{a=b=1}$。

\problem{54}{隐函数中带交叉分母的组合导数}{2015 年第 7 届，非数学类，一(2)，\src{07}}
$F(x+z/y,y+z/x)=0$ 确定 $z(x,y)$，$xy\ne0$，且 $xF_u+yF_v\ne0$。求 $xz_x+yz_y$（不含 $F$ 及其偏导）。
\sol 分别对 $x,y$ 求导，把含 $z_x,z_y$ 的项集中：
\[
z_x=\frac{y(zF_v-x^2F_u)}{xF_u+yF_v},\qquad
z_y=\frac{x(zF_u-y^2F_v)}{xF_u+yF_v}.
\]
分别乘以 $x,y$ 后相加，分子为 $(z-xy)(xF_u+yF_v)$，所以 $\ans{xz_x+yz_y=z-xy}$。求单个偏导时保留公共分母，目标组合往往会整体约掉。

\problem{55}{两层复合的混合二阶导数}{2023 年第 15 届，非数学 A/B 类，一(2)，共题，\src{15A}\src{15B}}
$z=f(x^2-y^2,xy)$，$f\in C^2$，求 $z_{xy}$。
\sol 令 $u=x^2-y^2,v=xy$，先得 $z_x=2xf_u+yf_v$。对 $y$ 求导：
\[z_{xy}=2x(-2yf_{uu}+xf_{uv})+f_v+y(-2yf_{vu}+xf_{vv}),\]
故 $\ans{z_{xy}=-4xyf_{uu}+2(x^2-y^2)f_{uv}+xyf_{vv}+f_v}$，所有外层偏导均在 $(u,v)$ 处取值。最后的 $f_v$ 来自系数 $y$ 的求导。

\problem{56}{指数内层也要二次求导}{2024 年第 16 届，非数学 A 类一(3)／B 类一(4)，共题，\src{16A}\src{16B}}
$z=f(xy,\ee^{x+y})$，$f\in C^2$，求 $z_{xy}$。
\sol 令 $u=xy,v=\ee^{x+y}$，$z_x=yf_u+vf_v$。再对 $y$ 求导并合并：
\[
\ans{z_{xy}=xyf_{uu}+(x+y)vf_{uv}+v^2f_{vv}+f_u+vf_v},\quad v=\ee^{x+y}.
\]
其中 $f_u$ 来自 $y$ 的导数，$vf_v$ 来自 $v$ 的导数；它们是二阶求导最常漏的项。

\problem{57}{先求隐函数值，再求偏导}{2025 年第 17 届，非数学 A/B 类，一(4)，共题，\src{17A}\src{17B}}
$xz=f(2x-y,xyz)$，$f(2,0)=3,f_u(2,0)=5$，$f\in C^1$。求 $z_x(1,0)$。
\sol 先由原式得 $z(1,0)=3$。固定 $y$ 求 $x$ 偏导：$z+xz_x=2f_u+f_v(yz+xyz_x)$。在 $(1,0)$，所有带 $y$ 的项消失，故 $3+z_x=10$，$\ans{z_x(1,0)=7}$。无需知道 $f_v(2,0)$。

\section{梯度、极值与约束真题}
\problem{58}{双约束最值可先消元}{2024 年第 16 届，非数学 B 类，二，\src{16B}}
在 $x-y+z=1$、$x^2+y^2=1$ 下，求 $x+2y+3z$ 的最大值及取值点。
\sol 消去 $z=1-x+y$，目标为 $3-2x+5y$。Cauchy 不等式给出 $-2x+5y\le\sqrt{29}\sqrt{x^2+y^2}=\sqrt{29}$，等号在 $(x,y)=(-2,5)/\sqrt{29}$。
故最大值 $\ans{3+\sqrt{29}}$，取值点 $\ans{(-2/\sqrt{29},5/\sqrt{29},1+7/\sqrt{29})}$。若用乘子法，消元后的方程是 $(-2,5)=2\lambda(x,y)$，同样得到此点。

\problem{59}{方向导数信息确定曲面，再求面积}{2025 年第 17 届，非数学 B 类，二，\src{17B}}
$f(x,y)=ax^2+by^2+1$。在 $(4,3)$ 处沿指向原点的方向取得最大方向导数 10。求 $a,b$，并求曲面 $z=f(x,y)$ 位于 $z=0$ 上方部分的面积。
\sol 指向原点的单位方向为 $(-4/5,-3/5)$，最大值为梯度模长，所以 $\grad f(4,3)=10(-4/5,-3/5)=(-8,-6)$。而梯度是 $(8a,6b)$，得 $\ans{a=b=-1}$。
曲面 $z=1-x^2-y^2$，投影为单位圆盘，面积
\[S=\iint_{x^2+y^2\le1}\sqrt{1+4x^2+4y^2}\dd x\dd y
=2\pi\int_0^1r\sqrt{1+4r^2}\dd r=\ans{\frac\pi6(5\sqrt5-1)}.\]
用方向向量前必须单位化，直接把 $(-4,-3)$ 当方向余弦会把梯度放大 5 倍。

\problem{60}{隐函数两个分支的局部极值}{2025 年第 17 届，非数学 B 类，三，\src{17B}}
由 $2x^2+2y^2+z^2+8xz-z+8=0$ 确定的局部隐函数 $z(x,y)$，求极值。
\sol 记左端为 $F$。驻点要求 $F_x=4x+8z=0$、$F_y=4y=0$，所以 $x=-2z,y=0$。代回得 $7z^2+z-8=0$，即 $z=1,-8/7$。
在 $(-2,0,1)$，$F_z=2z+8x-1=-15\ne0$；驻点处 $z_{xx}=z_{yy}=-4/F_z=4/15$、$z_{xy}=0$，故为极小值 $\ans1$。
在 $(16/7,0,-8/7)$，$F_z=15$，Hessian 负定，故为极大值 $\ans{-8/7}$。两者属于不同局部支，不能因“极小值比另一个极大值大”而认为矛盾。

\problem{61}{沿线段把多元问题化为一元}{2018 年第 10 届，非数学类，五，\src{10}}
$f$ 在域内可微，$f_x^2+f_y^2\le M$。若线段 $AB$ 全在域内，证明 $|f(B)-f(A)|\le\sqrt M\,|AB|$。
\sol 令 $\phi(t)=f(A+t(B-A))$，$0\le t\le1$。链式法则及 Cauchy 给出 $|\phi'(t)|=|\grad f\cdot(B-A)|\le\sqrt M|AB|$。用 Lagrange 或积分即得。
\trap{线段必须留在定义域中。任意连通域不必包含连接两点的直线段；此时只能沿域内路径估计长度。}

\chapter{重积分与变量替换}
\section{二重、三重积分规则}
\ruleitem{积分意义}{连续函数在闭有界区域可积，且 Fubini 定理允许按合适的先后次序累次积分。区域边界对内部变量常要分段；非矩形上不能遗漏上下限的依赖。}
\ruleitem{换元 Jacobian}{若 $T:(u,v)\mapsto(x,y)$ 可逆且光滑，则
\[
\iint_D f(x,y)\dd x\dd y
=\iint_{D'} f(x(u,v),y(u,v))
\left|\frac{\partial(x,y)}{\partial(u,v)}\right|\dd u\dd v.
\]
面积元取绝对值，不保留方向。变换须在积分区域上适当一一覆盖；若多重覆盖需计重数或切区。}
\ruleitem{极坐标}{$x=r\cos\theta,y=r\sin\theta$，$\dd x\dd y=r\dd r\dd\theta$。圆盘 $0\le r\le R,0\le\theta<2\pi$；环形扇区与含 $x^2+y^2$ 的被积函数优先考虑。}
\ruleitem{柱坐标}{$x=r\cos\theta,y=r\sin\theta,z=z$，$\dd V=r\dd r\dd\theta\dd z$。}
\ruleitem{球坐标}{约定 $\phi$ 为自正 $z$ 轴量起的极角，$0\le\phi\le\pi$，$\theta$ 为方位角：
\[
x=\rho\sin\phi\cos\theta,\ y=\rho\sin\phi\sin\theta,\ z=\rho\cos\phi,\quad
\dd V=\rho^2\sin\phi\dd\rho\dd\phi\dd\theta.
\]
球 $0\le\rho\le R$；上半球 $0\le\phi\le\pi/2$；圆锥 $z\ge\sqrt{x^2+y^2}$ 对应 $0\le\phi\le\pi/4$。}
\ruleitem{平移椭球为球}{若 $u=x-a,v=y-b,w=c(z-d)$，则 $\dd x\dd y\dd z=|c|^{-1}\dd u\dd v\dd w$。先移到球心，再缩放；线性换元总是要乘 Jacobian。}
\ruleitem{对称积分}{关于 $x$ 轴、$y$ 轴或原点对称的区域，奇被积函数积分为 0。球 $B_R$：$\iiint x^2\dd V=\iiint y^2\dd V=\iiint z^2\dd V=\frac13\iiint r^2\dd V=4\pi R^5/15$；球面 $S_R$：$\iint_{S_R}x^2\dd S=4\pi R^4/3$，圆盘半径 $R$：$\iint x^2\dd A=\pi R^4/4$。}
\ruleitem{应用}{薄片质量 $\iint_D\rho\dd A$，体质量 $\iiint_\Omega\rho\dd V$；质心坐标=对应一阶矩/质量；转动惯量 $\iiint r_\perp^2\rho\dd V$。上下曲面 $z=u(x,y),v(x,y)$ 的体积为 $\iint_D|v-u|\dd A$。}

\section{区域识别与累次积分真题}
\problem{62}{曲线两侧符号面积}{2011 年第 3 届，非数学类，一(3)，\src{03}}
在 $D=[0,2]^2$ 上求 $\iint_D\sgn(xy-1)\dd x\dd y$。
\sol 负值区域 $xy<1$ 的面积分层为
\[A=\int_0^{1/2}2\dd x+\int_{1/2}^2\frac1x\dd x=1+2\ln2.\]
全域面积 4，所求为正面积减负面积 $4-2A=\ans{2-4\ln2}$。曲线 $y=1/x$ 要与上边界 $y=2$ 分段。
\problem{63}{两抛物面与切平面围成体积}{2015 年第 7 届，非数学类，一(3)，\src{07}}
求曲面 $z=x^2+y^2+1$ 在 $(1,-1,3)$ 处切平面与 $z=x^2+y^2$ 所围区域的体积。
\sol 梯度 $(2x,2y)$ 给切平面 $z=2x-2y-1$。两面高度差为
\[1-(x-1)^2-(y+1)^2.\]
投影是以 $(1,-1)$ 为圆心的单位圆。平移极坐标得
\[V=2\pi\int_0^1(1-r^2)r\dd r=\ans{\pi/2}.\]
先确定上下面和投影圆盘，再积分；体积不能对带符号高度差取绝对值而不检查区域。
\problem{64}{平移椭球后二次矩积分}{2016 年第 8 届，非数学类，三，\src{08}}
$\Omega:x^2+y^2+2z^2\le x+y+2z$，求 $\iiint_\Omega(x^2+y^2+z^2)\dd V$。
\sol 配方后中心为 $(1/2,1/2,1/2)$。令 $u=x-1/2,v=y-1/2,w=\sqrt2(z-1/2)$，变为单位球，$dV=du\,dv\,dw/\sqrt2$。中心对称消去一次项，故
\[
I=\frac1{\sqrt2}\iiint_{B_1}(u^2+v^2+w^2/2+3/4)\dd V
=\frac1{\sqrt2}\left(\frac52\frac{4\pi}{15}+\pi\right)
=\ans{\frac{5\sqrt2\pi}{6}}.
\]
配方平移不改变体积元；非均匀缩放才产生 $1/\sqrt2$。
\problem{65}{球坐标与圆锥区域}{2015 年第 7 届，非数学类，一(6)，\src{07}}
曲面 $z^2=x^2+y^2$ 与 $z=\sqrt{4-x^2-y^2}$ 围成区域 $V$，求 $\iiint_V z\dd V$。
\sol 上半球半径 2 与上锥面交于 $\phi=\pi/4$。在锥面上方、球面下方，
$0\le\rho\le2$、$0\le\phi\le\pi/4$、$0\le\theta<2\pi$。因此
\[
I=2\pi\int_0^2\rho^3\dd\rho\int_0^{\pi/4}\cos\phi\sin\phi\dd\phi
=\ans{2\pi}.
\]
球坐标的 $\phi$ 是从正 $z$ 轴量起；若改用从水平面量的仰角，上下限和 Jacobian 都要相应改变。
\problem{66}{交换积分次序揭示简洁区域}{2023 年第 15 届，非数学 B 类，一(5)，\src{15B}}
求 $\displaystyle I=\int_0^1\dd x\int_x^{\sqrt x}\frac{\cos y}{y}\dd y$。
\sol 区域约束等价于 $0\le y\le1$、$y^2\le x\le y$，所以
\[
I=\int_0^1(y-y^2)\frac{\cos y}y\dd y
=\int_0^1(1-y)\cos y\dd y
=\ans{1-\cos1}.
\]
原式虽然内层有 $1/y$，组合后的区域宽度 $y-y^2$ 抵消了端点奇异。
\problem{67}{圆盘上的小尺度极限}{2024 年第 16 届，非数学 A 类二／B 类五，共题，\src{16A}\src{16B}}
半径 $r$ 圆盘 $D_r$ 上求 $\lim_{r\to0^+}r^{-4}\iint_{D_r}(\ee^{x^2+y^2}-1)\dd A$。
\sol 极坐标直接积分：
\[
\iint_{D_r}(\ee^{x^2+y^2}-1)\dd A
=2\pi\int_0^r(\ee^{\rho^2}-1)\rho\dd\rho
=\pi(\ee^{r^2}-1-r^2)
=\frac\pi2r^4+o(r^4).
\]
故答案 $\ans{\pi/2}$。亦可先用 $\ee^u-1\sim u$，再用 $\iint_{D_r}(x^2+y^2)\dd A=\pi r^4/2$；尺度阶数不可漏掉面积元里的 $\rho$。

\chapter{曲线积分、曲面积分与三大公式}
\section{有向积分、参数化与公式条件}
\ruleitem{第一类曲线积分}{$\int_L f\dd s$ 不依赖走向。参数曲线 $\bm r(t)$ 的弧长元为 $|\bm r'(t)|\dd t$。}
\[\int_L f\dd s=\int_\alpha^\beta f(\bm r(t))|\bm r'(t)|\dd t.\]
分段光滑即可。
\ruleitem{第二类曲线积分}{$\int_L P\dd x+Q\dd y+R\dd z$ 依赖方向。参数化后为
$\int_\alpha^\beta(Px'+Qy'+Rz')\dd t$。反向取负；闭合曲线仍须说明正向。}
\ruleitem{第一类曲面积分}{$\iint_S f\dd S$ 无定向，参数曲面 $\bm r(u,v)$ 的面积元为 $|\bm r_u\times\bm r_v|\dd u\dd v$。图像 $z=g(x,y)$：$\dd S=\sqrt{1+g_x^2+g_y^2}\dd x\dd y$。}
\ruleitem{通量积分}{$\iint_S \bm F\cdot\bm n\dd S$ 依赖定向。参数化有向面积向量 $\bm r_u\times\bm r_v\dd u\dd v$；若取反向则积分变号。图像 $z=g$ 向上面积向量 $(-g_x,-g_y,1)\dd x\dd y$。}
\ruleitem{Green 公式}{平面正向简单闭曲线（逆时针）$C=\partial D$，$P,Q\in C^1$ 于包含 $D$ 的邻域：
\[
\oint_C P\dd x+Q\dd y=\iint_D(Q_x-P_y)\dd A.
\]
有孔区域边界方向为外边界逆时针、内边界顺时针。若场在区域内有奇点，不能直接套用。}
\ruleitem{保守场}{单连通域内 $Q_x=P_y$（三维 $\operatorname{curl}F=0$）推出存在势函数的常见充分条件；非单连通域只有必要而非充分。例 $(-y/(x^2+y^2),x/(x^2+y^2))$ 在去心平面旋度为零，但绕原点的闭路积分为 $2\pi$。}
\ruleitem{Gauss 散度定理}{闭合、分段光滑边界 $S=\partial\Omega$ 取外法向，$\bm F\in C^1$ 在包含区域的开集：
\[
\iint_{\partial\Omega}\bm F\cdot\bm n\dd S=\iiint_\Omega\operatorname{div}\bm F\dd V.
\]
若向量场在体内有奇点则不适用；内部不能凭空穿过奇点。}
\ruleitem{Stokes 公式}{定向曲面 $S$ 有边界 $C$，边界按右手法则与单位法向对应，向量场 $C^1$：
\[
\oint_C\bm F\cdot\dd\bm r=\iint_S(\operatorname{curl}\bm F)\cdot\bm n\dd S.
\]
取反定向两侧同时变号。$\operatorname{div}(\operatorname{curl}F)=0$，$\operatorname{curl}(\grad\phi)=0$。}

\section{通量、环流与几何真题}
\problem{68}{球面对线性组合积分}{2011 年第 3 届，非数学类，六，\src{03}}
球面单位球 $S$ 上求 $\iint_S f(ax+by+cz)\dd S$（$f$ 连续）。
\sol 旋转坐标，使 $(a,b,c)$ 转为 $z$ 轴。设 $A=\sqrt{a^2+b^2+c^2}$，单位球上 $z=t$ 的带状面积为 $2\pi\dd t$，所以
\[\ans{\iint_S f(ax+by+cz)\dd S=2\pi\int_{-1}^1f(At)\dd t}.\]
若 $A=0$，两边均为 $4\pi f(0)$。半径若非 1，面积带会多出 $R^2$ 因子。

\problem{69}{用散度定理直接找最小通量区域}{2013 年第 5 届，非数学类，五，\src{05}}
对闭合定向曲面（取外法向）与场 $\bm F=(x^3-x,2y^3-y,3z^3-z)$，求可能的最小通量。
\sol 散度为 $3(x^2+2y^2+3z^2-1)$。记 $q=x^2+2y^2+3z^2$。为使体积分最小，包含所有散度为负的点且排除散度为正的点，候选区域为 $q\le1$。
线性换元 $u=x,v=\sqrt2y,w=\sqrt3z$，Jacobian 为 $1/\sqrt6$。区域内被积函数为 $3(q-1)$，故
\[
\Phi_{\min}=\frac3{\sqrt6}\,4\pi\int_0^1(r^2-1)r^2\dd r
=\ans{-\frac{4\sqrt6\pi}{15}}.
\]
边界椭球上的通量由 Gauss 定理取得。题设若限定某类区域或允许多个边界组件，通常的体积分极小原理仍成立，但应明确区域边界光滑并避开病态边界。
\problem{70}{平面截椭圆柱的第一类面积分}{2023 年第 15 届，非数学 A 类，一(5)，\src{15A}}
$L$ 为平面 $y+z=5$ 与圆柱 $x^2+y^2=25$ 所围平面片，求 $\iint_L(x+y+z)\dd S$。
\sol 写成 $z=5-y$，$z_x=0,z_y=-1$，故 $\dd S=\sqrt2\dd x\dd y$，投影是半径 5 圆盘。被积函数化为 $x+5$；奇项 $x$ 积分为 0，故
\[\ans{\sqrt2\cdot5\cdot25\pi=125\sqrt2\pi}.\]

\problem{71}{绕奇点的环流不能用 Green}{2024 年第 16 届，非数学 A 类，四，\src{16A}}
令 $L:x^2+y^2=9$ 逆时针，求
$\oint_L[-y/(4x^2+y^2)\dd x+x/(4x^2+y^2)\dd y]$。
\sol 参数化 $(x,y)=(3\cos t,3\sin t)$。代入得被积参数 $1/(4\cos^2t+\sin^2t)$，所以
\[
I=4\int_0^{\pi/2}\frac{\dd t}{4\cos^2t+\sin^2t}
=4\int_0^\infty\frac{\dd u}{4+u^2}=\ans{\pi}.
\]
原点的分母为零，位于椭圆内部；区域不满足 Green 公式的光滑性条件。直接参数化既安全又简短。

\problem{72}{上半球通量与补底面}{2024 年第 16 届，非数学 A 类，四，\src{16A}}
上半球面 $S_R$ 取外法向，计算场 $\bm F=(x^2-x,y^2-y,z^2-z)$ 的通量。
\sol 补上圆盘 $z=0$ 成闭曲面。底面向下法向为 $-\bm k$，而场的 $z$ 分量为 0，底面通量为 0。散度为 $2(x+y+z)-3$。半球中 $x,y$ 奇项积分为零，$\iiint z\dd V=\pi R^4/4$，体积 $2\pi R^3/3$。故
\[
\Phi=\frac{\pi R^4}2-2\pi R^3=\ans{\pi R^3(R-4)/2}.
\]
$R>0$ 任意。底面法向方向不影响本题答案，但封闭计算必须说清楚取向。
\problem{73}{圆周上的二次多项式线积分}{2025 年第 17 届，非数学 A 类，一(5)，\src{17A}}
曲线 $L$ 是球面 $x^2+y^2+z^2=4$ 与平面 $x+y+z=3$ 的交线，求 $\int_L(x^2+2xy)\dd s$。
\sol 圆心为平面法向投影 $(1,1,1)$，即 $(1,1,1)$；半径 $\sqrt{4-3}=1$，所以周长 $2\pi$。写 $\bm r=(1,1,1)+\cos t\,\bm e_1+\sin t\,\bm e_2$，取
$\bm e_1=(1,-1,0)/\sqrt2$、$\bm e_2=(1,1,-2)/\sqrt6$。
一周期平均有 $\overline{x^2}=1+\frac12(1/2+1/6)=4/3$，$\overline{xy}=1+\frac12(-1/2+1/6)=5/6$。因此
\[\ans{\int_L(x^2+2xy)\dd s=2\pi(4/3+5/3)=6\pi}.\]
$\dd s=\dd t$（半径为 1）。“球面与平面的交线”先求圆心、半径，再用圆周平均值处理二次式。

\chapter{数项级数、幂级数与傅里叶级数}
\section{数项级数敛散判别}
级数 $\sum a_n$ 收敛的必要条件是 $a_n\to0$，但非充分。柯西准则要求任意长的尾部部分和趋于 0。经典级数：
\[
\sum_{n=0}^\infty q^n=\frac1{1-q}\ (|q|<1),\quad
\sum_{n=1}^\infty\frac1{n^p}\text{ 收敛}\iff p>1,\quad
\sum_{n=1}^\infty\frac1{n(\ln n)^p}\text{ 收敛}\iff p>1.
\]
调和级数发散。加有限项不改变敛散性。
\ruleitem{正项比较}{最终 $0\le a_n\le b_n$ 且 $\sum b_n$ 收敛，则 $\sum a_n$ 收敛；若 $a_n\ge b_n\ge0$ 且 $\sum b_n$ 发散，则 $\sum a_n$ 发散。若 $a_n/b_n\to c\in(0,\infty)$，则同敛散。}
\ruleitem{比值、根值}{$L=\lim|a_{n+1}/a_n|$ 或 $L=\limsup\sqrt[n]{|a_n|}$：$L<1$ 绝对收敛，$L>1$ 发散；$L=1$ 无结论。阶乘与指数常先用比值。}
\ruleitem{交错与误差}{若 $a_n\downarrow0$，则 $\sum(-1)^{n-1}a_n$ 收敛，截断误差 $|R_N|\le a_{N+1}$。单调条件与趋零缺一不可。绝对收敛必收敛；条件收敛重排可能改变和（提高）。}
\ruleitem{Dirichlet}{若 $A_N=\sum_{n=1}^N a_n$ 有界，$b_n$ 单调趋 0，则 $\sum a_nb_n$ 收敛。Abel：若 $\sum a_n$ 收敛、$b_n$ 单调有界，则 $\sum a_nb_n$ 收敛。先确认“部分和有界”或“单调”对应哪一列。}
\ruleitem{变换级数}{双重求和、积分与级数换序，须有非负项（Tonelli）或绝对/一致收敛等合法条件。条件收敛时不能任意交换次序。}

\section{幂级数与常用展开}
幂级数 $\sum c_n(x-a)^n$ 的收敛半径
\[
R=\frac1{\limsup_{n\to\infty}\sqrt[n]{|c_n|}},
\]
约定 $1/0=\infty,1/\infty=0$。若系数比值极限可用 $R=\lim|c_n/c_{n+1}|$。在 $|x-a|<R$ 内逐项求导、积分；边界 $|x-a|=R$ 逐点另判。
\begin{align*}
\frac1{1-x}&=\sum_{n=0}^\infty x^n &&(|x|<1),\\
\ln(1+x)&=\sum_{n=1}^\infty\frac{(-1)^{n-1}x^n}{n} &&(-1<x\le1),\\
\arctan x&=\sum_{n=0}^\infty\frac{(-1)^nx^{2n+1}}{2n+1} &&(|x|\le1),\\
\ee^x&=\sum_{n=0}^\infty\frac{x^n}{n!},\quad
\sin x=\sum_{n=0}^\infty\frac{(-1)^nx^{2n+1}}{(2n+1)!},\\
\cos x&=\sum_{n=0}^\infty\frac{(-1)^nx^{2n}}{(2n)!} &&(x\in\R),\\
(1+x)^\alpha&=\sum_{n=0}^\infty\binom\alpha n x^n &&(|x|<1).
\end{align*}
其中 $\binom\alpha n=\alpha(\alpha-1)\cdots(\alpha-n+1)/n!$。边界以必要条件 $c_n(x-a)^n\to0$ 起步，必要条件失败即发散；通过后仍需判别。

\section{傅里叶级数（提高）}
周期 $2L$ 可积函数在适当分段光滑条件下的傅里叶系数为
\[
a_0=\frac1L\int_{-L}^{L}f(x)\dd x,\quad
a_n=\frac1L\int_{-L}^{L}f(x)\cos\frac{n\pi x}L\dd x,\quad
b_n=\frac1L\int_{-L}^{L}f(x)\sin\frac{n\pi x}L\dd x.
\]
级数为 $a_0/2+\sum_{n\ge1}(a_n\cos(n\pi x/L)+b_n\sin(n\pi x/L))$。在连续点收敛到 $f(x)$，在跳跃点收敛到左右极限平均。分段 $C^1$ 是常见充分条件。连续周期函数由其全部傅里叶系数唯一确定（可用 Fejér 定理）。

\section{幂级数、递推与收敛真题}
\problem{74}{递推微分方程后求幂级数和}{2009 年第 1 届，非数学类，七，\src{01}}
$u_n'=u_n+x^{n-1}\ee^x$，$u_n(1)=\ee/n$。求 $\sum_{n=1}^\infty u_n(x)$ 及定义区间。
\sol 乘积分因子 $\ee^{-x}$，得 $(\ee^{-x}u_n)'=x^{n-1}$；用初值积分得 $u_n=\ee^x x^n/n$。因此
\[
\sum_{n=1}^\infty u_n(x)=-\ee^x\ln(1-x),\quad |x|<1.
\]
在 $x=-1$ 交错收敛，级数和由 Abel 定理得 $-\ee^{-1}\ln2$；$x=1$ 项为 $\ee/n$ 发散；$x<-1$ 通项不趋零。因此收敛集合 $\ans{[-1,1)}$。

\problem{75}{从几何级数求有理函数和}{2011 年第 3 届，非数学类，一(4)，\src{03}}
求 $S(x)=\sum_{n=1}^\infty(2n-1)x^{2n-2}/2^n$ 的和与定义域，并求 $\sum_{n=1}^\infty(2n-1)/2^{2n-1}$。
\sol 令 $q=x^2/2$，用 $\sum q^{n-1}=(1-q)^{-1}$ 及其导数，得
\[
S(x)=\frac12\sum_{n=1}^\infty(2n-1)q^{n-1}
=\frac{1+q}{2(1-q)^2}=\frac{2+x^2}{(2-x^2)^2}.
\]
根值检验得 $|x|<\sqrt2$，端点通项不趋 0。所求数值为 $S(1/\sqrt2)=\ans{10/9}$。

\problem{76}{二阶导数控制函数值级数}{2013 年第 5 届，非数学类，二，\src{05}}
$f$ 在 0 二阶可导，且 $\lim_{x\to0}f(x)/x=0$。证明 $\sum_{n=1}^\infty|f(1/n)|$ 收敛。
\sol 极限先给出 $f(0)=0$。由差商定义 $f'(0)=0$。二阶 Taylor--Peano 式给
$f(x)=\frac12f''(0)x^2+o(x^2)=O(x^2)$，所以 $|f(1/n)|\le C/n^2$ 对充分大的 $n$ 成立。比较 $p$ 级数，结论成立。这里二阶可导足以有 Peano 展开，不需要假设 $f''$ 连续。

\problem{77}{调和数与非负双重求和}{2013 年第 5 届，非数学类，七，\src{05}}
求 $\displaystyle\sum_{n=1}^\infty\frac{H_n}{(n+1)(n+2)}$，其中 $H_n=\sum_{k=1}^n1/k$。
\sol 项均非负，可先按有限矩形求和，再由单调极限交换次序：
\[
\sum_{n=1}^\infty\sum_{k=1}^n\frac1{k(n+1)(n+2)}
=\sum_{k=1}^\infty\frac1k\sum_{n=k}^\infty\left(\frac1{n+1}-\frac1{n+2}\right)
=\sum_{k=1}^\infty\frac1{k(k+1)}=\ans1.
\]
非负项交换不依赖先验收敛，若换成带符号项则须另证绝对/一致收敛。

\problem{78}{跳跃点傅里叶和}{2015 年第 7 届，非数学类，一(4)，\src{07}}
将 $f(x)=3$（$-5\le x<0$）、$f(x)=0$（$0\le x<5$）作周期 10 延拓，求其傅里叶级数在 $x=0$ 的和。
\sol 周期延拓在 0 的左极限为 3、右极限为 0。Dirichlet 收敛定理给跳跃平均值 $\ans{3/2}$。级数的和并非函数在跳点所定义的 $f(0)=0$；边界约定不改变跳跃平均。

\problem{79}{阶乘分母的整系数拆分}{2015 年第 7 届，非数学类，四，\src{07}}
求 $\displaystyle\sum_{n=0}^\infty\frac{n^3+2}{(n+1)!}(x-1)^n$ 及收敛半径。
\sol 令 $t=x-1,m=n+1$。因为 $(n^3+2)=(m-1)^3+2=m^3-3m^2+3m+1$，且
$\sum m^3t^m/m!=(t^3+3t^2+t)\ee^t$、$\sum m^2t^m/m!=(t^2+t)\ee^t$、$\sum mt^m/m!=t\ee^t$，整理后
\[
\ans{\sum_{n=0}^\infty\frac{n^3+2}{(n+1)!}(x-1)^n
=\frac{(t^3+t+1)\ee^t-1}{t}}.
\]
$t=0$ 取连续延拓，值为 2。阶乘使半径 $\ans{R=\infty}$。可用首两项复核封闭式：常数项 2，$t$ 项 $3/2$。

\problem{80}{对幂级数求导再积分}{2023 年第 15 届，非数学 B 类，四，\src{15B}}
求 $S(x)=\displaystyle\sum_{n=1}^\infty\frac{(-1)^{n-1}x^{2n}}{n(2n-1)}$ 的收敛域与和。
\sol 半径为 1。对 $|x|<1$ 求导两次：
\[
S''(x)=2\sum_{n=1}^\infty(-1)^{n-1}x^{2n-2}=\frac2{1+x^2}.
\]
$S(0)=S'(0)=0$，积分得 $S'(x)=2\arctan x$、$S(x)=2x\arctan x-\ln(1+x^2)$。在端点原级数绝对收敛，因为通项绝对值为 $O(n^{-2})$；端点由连续延拓，故 $\ans{x\in[-1,1]}$。

\problem{81}{部分分式决定 Taylor 半径与端点}{2024 年第 16 届，非数学 B 类，三，\src{16B}}
求 $f(x)=1/[(x-1)(x+3)]$ 在 $x=2$ 的 Taylor 级数及收敛区间。
\sol 令 $h=x-2$，
\[
f(x)=\frac14\left(\frac1{1+h}-\frac1{5+h}\right)
=\frac14\sum_{n=0}^\infty(-1)^n(1-5^{-n-1})h^n.
\]
最近奇点距中心 1，故半径为 1，区间 $(1,3)$。在 $h=1$ 通项不趋零；在 $h=-1$ 原函数有极点且通项不趋零，两个端点都发散。

\problem{82}{递推差分望远镜成正项级数和}{2023 年第 15 届，非数学 A 类，六，\src{15A}}
$x_0=1/3$，$x_{n+1}=x_n^2/(1-x_n+x_n^2)$。求 $\sum_{n=0}^\infty x_n$。
\sol 对 $0<x<1$，有 $0<x^2/(1-x+x^2)<x$，故 $x_n\downarrow0$。直接通分可得
\[
\frac{x_n}{1-x_n}-\frac{x_{n+1}}{1-x_{n+1}}=x_n.
\]
求和望远镜：$\sum_{n=0}^N x_n=\frac{x_0}{1-x_0}-\frac{x_{N+1}}{1-x_{N+1}}\to\ans{1/2}$。

\problem{83}{两个不可公度周期迫使 Fourier 系数消失}{2016 年第 8 届，非数学类，六，\src{08}}
$f$ 可微，且 $f(x+2)=f(x+\sqrt3)=f(x)$。用 Fourier 方法证明 $f$ 为常数。
\sol 将 $f$ 看作周期 2 的连续函数，系数
$c_n=\frac12\int_{-1}^1f(x)\ee^{-in\pi x}\dd x$。由周期 2 的函数平移不变性，对 $\sqrt3$ 作换元可得
$c_n=\ee^{in\pi\sqrt3}c_n$。若 $n\ne0$，$n\sqrt3$ 非整数，故相位不为 1，$c_n=0$。所有非零频率系数消失，傅里叶唯一性（Fejér 定理）推出 $f=c_0$。

\problem{84}{积分定义数列的敛散、渐近与精确和}{2020 年第 12 届，非数学类，七，\src{12}}
$u_n=\int_0^1(1+t^4)^{-n}\dd t$。讨论 $\sum(-1)^{n-1}u_n$、$\sum u_n$、$\sum u_n/n^p$，并求 $p=1$ 时的和。
\sol 被积函数随 $n$ 严格下降，且 $u_n\to0$（由控制收敛，$t>0$ 时点态趋零），所以交错级数收敛。对 $0\le t\le n^{-1/4}$，$(1+t^4)^n\le(1+1/n)^n<\ee$，故 $u_n\ge\ee^{-1}n^{-1/4}$；于是 $\sum u_n$ 发散。
对 $p=1$，非负项 Tonelli 允许求和与积分交换：
\[
\sum_{n=1}^\infty\frac{u_n}n
=\int_0^1\sum_{n=1}^\infty\frac{(1+t^4)^{-n}}n\dd t
=\int_0^1\ln(1+t^{-4})\dd t
=\ln2+4u_1.
\]
最后一步将对数写成 $\ln(1+t^4)-4\ln t$，对第一项分部。又
$u_1=\int_0^1\frac{\dd t}{1+t^4}
=\frac{\sqrt2}{8}[\pi+2\ln(1+\sqrt2)]$，故
$\ans{\sum u_n/n=\ln2+\frac{\sqrt2}{2}[\pi+2\ln(1+\sqrt2)]}$。
对 $p>1$，$0<u_n\le1$ 且 $u_n/n^p\le1/n^p$，故收敛。$p<1$ 时由于 $u_n\ge c n^{-1/4}$，若 $p<3/4$ 级数发散；$3/4\le p<1$ 需要更精确渐近，不能由上述粗下界断言。实际上令 $s=n^{1/4}t$ 并用 Laplace 极限得 $u_n\sim n^{-1/4}\int_0^\infty\ee^{-s^4}\dd s=\Gamma(1/4)/(4n^{1/4})$，故 $\sum u_n/n^p$ 收敛当且仅当 $p>3/4$。在 $p=3/4$ 比较调和级数发散。

\appendix
\chapter{考场方法速查与复习安排}
\section{题型信号 → 优先方法}
\begin{longtable}{p{.29\textwidth}p{.59\textwidth}}
\toprule 题面信号 & 先试的动作与检查点\\\midrule\endhead
$n\to\infty$，多项乘积 & 望远镜、取对数、等价首项、单调有界；对数取值前检查正性\\
$1^\infty$、$(1+u_n)^{v_n}$ & 取 $v_n\ln(1+u_n)$，确认乘积极限\\
差分 $g(n+1)-g(n)$ & 中值定理或 Stolz；Stolz 要核实分母单调趋无穷\\
正弦在等距点求和 & 识别黎曼和，先确认归一化步长\\
分子中连续多项相减 & 泰勒展开到抵消后的首个非零项；每个等价量要对应同一趋近变量\\
零点存在或根数 & 介值定存在，严格单调定唯一，Rolle 反复使用定上界\\
两点函数值差 & Lagrange；两个区间内各有导数点常先造一个中间分点\\
$f'+pf$ 与所求等于零 & 乘 $\ee^{\int p}$ 构造导数\\
积分式等于函数 & 先设积分为常数，或求导消去积分，再回代定常数\\
反常积分的振荡因子 & Dirichlet / Abel；绝对收敛单独判，瑕点分段\\
$x^2+y^2$、圆盘、径向函数 & 极坐标；面积元别漏 $r$\\
球/锥/球壳、$x^2+y^2+z^2$ & 球坐标；明确极角是从轴还是从平面量起\\
椭球二次式 & 配方、平移、缩放成球；列出 Jacobian\\
曲线 $P\dd x+Q\dd y$ 闭合 & Green 前查内部无奇点及方向\\
闭曲面通量 & Gauss 前查闭合、外法向、场在区域内光滑\\
边界曲线环流与旋度 & Stokes 检查右手定向；可换为容易计算的跨曲面\\
高阶复合偏导 & 先命名内层变量，列一阶链式式子，再逐项求导\\
隐函数偏导 & 先由原式找所求点的函数值，再对指定变量求导并代点\\
约束最值且约束简单 & 优先消元；Lagrange 提供候选，最后检查边界与奇点\\
$n!$ 或 $a^n$ 系数 & 比值法求收敛半径，幂级数封闭和先尝试系数拆分\\
数列项定义成积分 & 积分估计单调与趋零；换序须非负或绝对收敛\\
\bottomrule
\end{longtable}

\section{考场十条检查}
\begin{enumerate}
\item 极限题先写趋近方向，$x\to0^+$ 与双侧不能混。
\item 泰勒展开必须写清阶数，若首项抵消就继续展开。
\item 中值定理逐条核对闭区间连续、开区间可导、端点关系。
\item 罗必达只针对适用的不定式；每次使用后重新核验。
\item 不定积分不要漏绝对值、定义区间与 $C$；定积分换元同步换限。
\item 反常积分在每个瑕点分别取极限，禁止把发散项相减。
\item 方向导数用单位向量；切线、曲面面积、坐标变换的系数分别核对。
\item 场积分先判断第一类/第二类、是否有向、曲线或曲面方向。
\item 幂级数端点与内部半径分开判，收敛域通常写区间端点。
\item 计算完代回原题；极值候选检查约束，隐函数候选检查非零条件。
\end{enumerate}

\section{四周复习路线}
\begin{longtable}{p{.15\textwidth}p{.37\textwidth}p{.39\textwidth}}
\toprule 周次 & 学习内容 & 每周产出\\\midrule\endhead
第 1 周 & 极限、连续、导数、泰勒 & 整理一张展开表；完成真题 01—21；能解释每个极限展开到几阶\\
第 2 周 & 中值定理、积分、反常积分、ODE & 完成 22—47；独立写出两个证明构造、三种积分判敛证明\\
第 3 周 & 空间几何、多元微分、重积分 & 完成 48—67；练习隐函数偏导和坐标区域草图\\
第 4 周 & 线面积分、级数、混合套题 & 完成 68—84；至少两次限时模拟，按错因补漏\\
\bottomrule
\end{longtable}
时间更短时，按章节序号分为基础计算、综合计算、证明与提高三层；先保证每种题型至少完成一道，再集中补弱项。复盘需记“错误原因—正确触发信号—下次检查项”，不要只誊答案。

\section{一页式公式回看}
\[
\begin{array}{ll}
\text{等价首项：}&\sin x\sim x,\ \ln(1+x)\sim x,\ \ee^x-1\sim x,\ 1-\cos x\sim x^2/2.\\
\text{中值：}&f(b)-f(a)=f'(\xi)(b-a).\\
\text{Taylor：}&f(a+h)=\sum_{k=0}^n f^{(k)}(a)h^k/k!+o(h^n).\\
\text{积分分部：}&\int u\dd v=uv-\int v\dd u.\\
\text{线性 ODE：}&y'+Py=Q,\quad y=\ee^{-\int P}(\int Q\ee^{\int P}+C).\\
\text{方向导数：}&D_{\bm e}f=\grad f\cdot\bm e,\quad|\bm e|=1.\\
\text{二元曲面面积：}&\dd S=\sqrt{1+z_x^2+z_y^2}\dd x\dd y.\\
\text{极坐标：}&\dd A=r\dd r\dd\theta.\\
\text{球坐标：}&\dd V=\rho^2\sin\phi\dd\rho\dd\phi\dd\theta.\\
\text{Green：}&\oint(P\dd x+Q\dd y)=\iint(Q_x-P_y)\dd A.\\
\text{Gauss：}&\iint_{\partial\Omega}\bm F\cdot\bm n\dd S=\iiint_\Omega\div\bm F\dd V.\\
\text{Stokes：}&\oint_{\partial S}\bm F\cdot\dd\bm r=\iint_S\curl\bm F\cdot\bm n\dd S.\\
\text{幂级数：}&R^{-1}=\limsup|c_n|^{1/n};\quad\text{端点各自判定。}
\end{array}
\]
此表供考前回看，完整条件仍以各章定理说明为准。

\chapter{历年试卷来源与题号核对}
题目编号和年度以下载到并核查的试卷/答案页为准。S01—S14 为学校数学资料页打包转载的往届卷；S15—S17 为数学竞赛资料网站发布的近年非数学类卷及参考答案。后者作为可核查试卷副本使用，不称为竞赛官方发布。官方报名、赛制、当年考试安排和考纲应另查当年竞赛通知。

\section{来源清单}
\begin{longtable}{p{.13\textwidth}p{.42\textwidth}p{.36\textwidth}}
\toprule 编号 & 试卷 & 可核查入口\\\midrule\endhead
\srcentry{01} & 第 1 届，2009，非数学类预赛 & \href{https://sxjm.cdcas.edu.cn/rootfiles/2024/10/26/1728353777295545-1728353777378289.zip}{试题与解答 ZIP}\\
\srcentry{02} & 第 2 届，2010，非数学类初赛 & \href{https://sxjm.cdcas.edu.cn/rootfiles/2024/10/26/1728353777561057-1728353777573298.zip}{试题与解答 ZIP}\\
\srcentry{03} & 第 3 届，2011，非数学类初赛 & \href{https://sxjm.cdcas.edu.cn/rootfiles/2024/10/26/1728353777671018-1728353777680210.zip}{试题与解答 ZIP}\\
\srcentry{04} & 第 4 届，2012，非数学类初赛 & \href{https://sxjm.cdcas.edu.cn/rootfiles/2024/10/26/1728353777983692-1728353777993742.zip}{试题与解答 ZIP}\\
\srcentry{05} & 第 5 届，2013，非数学类初赛 & \href{https://sxjm.cdcas.edu.cn/rootfiles/2024/10/26/1728353778099713-1728353778114189.zip}{试题与解答 ZIP}\\
\srcentry{06} & 第 6 届，2014，非数学类预赛 & \href{https://sxjm.cdcas.edu.cn/rootfiles/2024/10/26/1728353778584723-1728353778599794.zip}{试题与解答 ZIP}\\
\srcentry{07} & 第 7 届，2015，非数学类预赛 & \href{https://sxjm.cdcas.edu.cn/rootfiles/2024/10/26/1728353778773978-1728353778794206.zip}{试题与解答 ZIP}\\
\srcentry{08} & 第 8 届，2016，非数学类预赛 & \href{https://sxjm.cdcas.edu.cn/rootfiles/2024/10/26/1728353778944297-1728353778953674.zip}{试题与解答 ZIP}\\
\srcentry{09} & 第 9 届，2017，非数学类预赛 & \href{https://sxjm.cdcas.edu.cn/rootfiles/2024/10/26/1728353779004697-1728353779017541.zip}{试题与解答 ZIP}\\
\srcentry{10} & 第 10 届，2018，非数学类预赛 & \href{https://sxjm.cdcas.edu.cn/rootfiles/2024/10/26/1728353779066362-1728353779078094.zip}{试题与解答 ZIP}\\
\srcentry{11} & 第 11 届，2019，非数学类初赛 & \href{https://sxjm.cdcas.edu.cn/rootfiles/2024/10/26/1728353779353603-1728353779366674.zip}{试题与解答 ZIP}\\
\srcentry{12} & 第 12 届，2020，非数学类初赛 & \href{https://sxjm.cdcas.edu.cn/rootfiles/2024/10/26/1728353779422669-1728353779435419.zip}{试题与解答 ZIP}\\
\srcentry{13} & 第 13 届，初赛补赛 & \href{https://sxjm.cdcas.edu.cn/rootfiles/2024/10/26/1728353779677414-1728353779689995.zip}{补赛试题与答案 ZIP}\\
\srcentry{14} & 第 14 届，初赛补赛（二） & \href{https://sxjm.cdcas.edu.cn/rootfiles/2024/10/26/1728353779744091-1728353779757371.zip}{补赛二试题 ZIP}\\
\srcentry{15A} & 第 15 届，2023，非数学 A 类 & \href{https://www.cmathc.org.cn/static/upload/file/jsst/%E7%AC%AC%E5%8D%81%E4%BA%94%E5%B1%8A%E5%85%A8%E5%9B%BD%E5%A4%A7%E5%AD%A6%E7%94%9F%E6%95%B0%E5%AD%A6%E7%AB%9E%E8%B5%9B%E5%88%9D%E8%B5%9B%EF%BC%88%E9%9D%9E%E6%95%B0%E5%AD%A6A%E7%B1%BB%EF%BC%89%E8%AF%95%E9%A2%98%E5%8F%8A%E5%8F%82%E8%80%83%E7%AD%94%E6%A1%88.pdf}{试题与参考答案 PDF}\\
\srcentry{15B} & 第 15 届，2023，非数学 B 类 & \href{https://www.cmathc.org.cn/static/upload/file/jsst/%E7%AC%AC%E5%8D%81%E4%BA%94%E5%B1%8A%E5%85%A8%E5%9B%BD%E5%A4%A7%E5%AD%A6%E7%94%9F%E6%95%B0%E5%AD%A6%E7%AB%9E%E8%B5%9B%E5%88%9D%E8%B5%9B%EF%BC%88%E9%9D%9E%E6%95%B0%E5%AD%A6B%E7%B1%BB%EF%BC%89%E8%AF%95%E9%A2%98%E5%8F%8A%E5%8F%82%E8%80%83%E7%AD%94%E6%A1%88.pdf}{试题与参考答案 PDF}\\
\srcentry{16A} & 第 16 届，2024，非数学 A 类 & \href{https://www.cmathc.org.cn/static/upload/file/jsst/%E7%AC%AC%E5%8D%81%E5%85%AD%E5%B1%8A%E5%85%A8%E5%9B%BD%E5%A4%A7%E5%AD%A6%E7%94%9F%E6%95%B0%E5%AD%A6%E7%AB%9E%E8%B5%9B%E5%88%9D%E8%B5%9B%EF%BC%88%E9%9D%9E%E6%95%B0%E5%AD%A6A%E7%B1%BB%EF%BC%89%E8%AF%95%E9%A2%98%E5%8F%8A%E7%AD%94%E6%A1%88%E8%A7%A3%E6%9E%90.pdf}{试题与答案解析 PDF}\\
\srcentry{16B} & 第 16 届，2024，非数学 B 类 & \href{https://www.cmathc.org.cn/static/upload/file/jsst/%E7%AC%AC%E5%8D%81%E5%85%AD%E5%B1%8A%E5%85%A8%E5%9B%BD%E5%A4%A7%E5%AD%A6%E7%94%9F%E6%95%B0%E5%AD%A6%E7%AB%9E%E8%B5%9B%E5%88%9D%E8%B5%9B%EF%BC%88%E9%9D%9E%E6%95%B0%E5%AD%A6B%E7%B1%BB%EF%BC%89%E8%AF%95%E9%A2%98%E5%8F%8A%E7%AD%94%E6%A1%88%E8%A7%A3%E6%9E%90.pdf}{试题与答案解析 PDF}\\
\srcentry{17A} & 第 17 届，2025，非数学 A 类 & \href{https://www.cmathc.org.cn/static/upload/file/jsst/%E7%AC%AC%E5%8D%81%E4%B8%83%E5%B1%8A%E5%85%A8%E5%9B%BD%E5%A4%A7%E5%AD%A6%E7%94%9F%E6%95%B0%E5%AD%A6%E7%AB%9E%E8%B5%9B%E5%88%9D%E8%B5%9B%E8%AF%95%E9%A2%98%E5%8F%8A%E5%8F%82%E8%80%83%E7%AD%94%E6%A1%88%EF%BC%88%E9%9D%9E%E6%95%B0%E5%AD%A6A%E7%B1%BB%EF%BC%89.pdf}{试题与参考答案 PDF}\\
\srcentry{17B} & 第 17 届，2025，非数学 B 类 & \href{https://www.cmathc.org.cn/static/upload/file/jsst/%E7%AC%AC%E5%8D%81%E4%B8%83%E5%B1%8A%E5%85%A8%E5%9B%BD%E5%A4%A7%E5%AD%A6%E7%94%9F%E6%95%B0%E5%AD%A6%E7%AB%9E%E8%B5%9B%E5%88%9D%E8%B5%9B%E8%AF%95%E9%A2%98%E5%8F%8A%E5%8F%82%E8%80%83%E7%AD%94%E6%A1%88%EF%BC%88%E9%9D%9E%E6%95%B0%E5%AD%A6B%E7%B1%BB%EF%BC%89.pdf}{试题与参考答案 PDF}\\
\bottomrule
\end{longtable}

\section{来源页面与使用说明}
\begin{itemize}
\item 成都文理学院数学教学相关资料页集中列有第 1—14 届（含补赛）下载包：\url{https://sxjm.cdcas.edu.cn/wjxz/sxjsst/C700502index_1.htm}。表中每条链接进入该页所列压缩包的直链。
\item 近年 2024、2025 试卷及答案可由数学竞赛资料网站的届次页核查：\url{https://www.cmathc.org.cn/csst/323.html}、\url{https://www.cmathc.org.cn/csst/359.html}；表内链接直接到对应 A/B PDF。
\item 竞赛当年的赛制、参赛类别、指定教材或考试范围，需以当年竞赛通知为准。西安交通大学 2023 届组织信息页可作 A/B 分类历史背景参考：\url{https://2023math.icrp.xjtu.edu.cn/}，不代表 2026 年通知。
\end{itemize}
题干在正文中均作简洁转述，仅保留解题所需条件；完整题目请点开原卷核对。参考答案供定位题意，不作为本资料推导的唯一依据。尤其是旧版扫描卷的符号，需结合原页图像辨读；正文已对易混符号与计算独立复核。

\end{document}
